% Options for packages loaded elsewhere
%DIF LATEXDIFF DIFFERENCE FILE
%DIF DEL Previously submitted (c952aa5)
%DIF ADD Minimal revision
%DIF LATEXDIFF DIFFERENCE FILE
%DIF DEL /tmp/phd029-submitted-notables.tex   Fri Sep 11 12:34:13 2026
%DIF ADD /tmp/phd029-current-notables.tex     Fri Sep 11 12:34:13 2026
\PassOptionsToPackage{unicode}{hyperref}
\PassOptionsToPackage{hyphens}{url}
%
\documentclass[
  12pt,
]{article}
\usepackage{amsmath,amssymb}
\usepackage{setspace}
\usepackage{iftex}
\ifPDFTeX
  \usepackage[T1]{fontenc}
  \usepackage[utf8]{inputenc}
  \usepackage{textcomp} % provide euro and other symbols
\else % if luatex or xetex
  \usepackage{unicode-math} % this also loads fontspec
  \defaultfontfeatures{Scale=MatchLowercase}
  \defaultfontfeatures[\rmfamily]{Ligatures=TeX,Scale=1}
\fi
\usepackage{lmodern}
\ifPDFTeX\else
  % xetex/luatex font selection
  \setmainfont[]{TeX Gyre Pagella}
\fi
% Use upquote if available, for straight quotes in verbatim environments
\IfFileExists{upquote.sty}{\usepackage{upquote}}{}
\IfFileExists{microtype.sty}{% use microtype if available
  \usepackage[]{microtype}
  \UseMicrotypeSet[protrusion]{basicmath} % disable protrusion for tt fonts
}{}
\makeatletter
\@ifundefined{KOMAClassName}{% if non-KOMA class
  \IfFileExists{parskip.sty}{%
    \usepackage{parskip}
  }{% else
    \setlength{\parindent}{0pt}
    \setlength{\parskip}{6pt plus 2pt minus 1pt}}
}{% if KOMA class
  \KOMAoptions{parskip=half}}
\makeatother
\usepackage{xcolor}
%DIF 42a42
\usepackage{graphicx} %DIF >
%DIF -------
\usepackage[margin=1in]{geometry}
\usepackage{longtable,booktabs,array}
\usepackage{calc} % for calculating minipage widths
% Correct order of tables after \paragraph or \subparagraph
\usepackage{etoolbox}
\makeatletter
\patchcmd\longtable{\par}{\if@noskipsec\mbox{}\fi\par}{}{}
\makeatother
% Allow footnotes in longtable head/foot
\IfFileExists{footnotehyper.sty}{\usepackage{footnotehyper}}{\usepackage{footnote}}
\makesavenoteenv{longtable}
\setlength{\emergencystretch}{3em} % prevent overfull lines
\providecommand{\tightlist}{%
  \setlength{\itemsep}{0pt}\setlength{\parskip}{0pt}}
\setcounter{secnumdepth}{-\maxdimen} % remove section numbering
\usepackage[natbibapa]{apacite}
\usepackage{lineno}\linenumbers
\ifLuaTeX
  \usepackage{selnolig}  % disable illegal ligatures
\fi
\IfFileExists{bookmark.sty}{\usepackage{bookmark}}{\usepackage{hyperref}}
\IfFileExists{xurl.sty}{\usepackage{xurl}}{} % add URL line breaks if available
\urlstyle{same}
\hypersetup{
  pdftitle={Cross-Cultural Validation and Ideological Fairness of a Historical Perspective Taking Instrument: Evidence from Czech Secondary Students},
  pdfauthor={Juda Kaleta},
  hidelinks,
  pdfcreator={LaTeX via pandoc}}

\title{Cross-Cultural Validation and Ideological Fairness of a Historical Perspective Taking Instrument: Evidence from Czech Secondary Students}
\author{Juda Kaleta}
\date{2026}
%DIF PREAMBLE EXTENSION ADDED BY LATEXDIFF
%DIF UNDERLINE PREAMBLE %DIF PREAMBLE
\RequirePackage[normalem]{ulem} %DIF PREAMBLE
\RequirePackage{color}\definecolor{RED}{rgb}{1,0,0}\definecolor{BLUE}{rgb}{0,0,1} %DIF PREAMBLE
\providecommand{\DIFaddtex}[1]{{\protect\color{blue}\uwave{#1}}} %DIF PREAMBLE
\providecommand{\DIFdeltex}[1]{{\protect\color{red}\sout{#1}}} %DIF PREAMBLE
%DIF SAFE PREAMBLE %DIF PREAMBLE
\providecommand{\DIFaddbegin}{} %DIF PREAMBLE
\providecommand{\DIFaddend}{} %DIF PREAMBLE
\providecommand{\DIFdelbegin}{} %DIF PREAMBLE
\providecommand{\DIFdelend}{} %DIF PREAMBLE
\providecommand{\DIFmodbegin}{} %DIF PREAMBLE
\providecommand{\DIFmodend}{} %DIF PREAMBLE
%DIF FLOATSAFE PREAMBLE %DIF PREAMBLE
\providecommand{\DIFaddFL}[1]{\DIFadd{#1}} %DIF PREAMBLE
\providecommand{\DIFdelFL}[1]{\DIFdel{#1}} %DIF PREAMBLE
\providecommand{\DIFaddbeginFL}{} %DIF PREAMBLE
\providecommand{\DIFaddendFL}{} %DIF PREAMBLE
\providecommand{\DIFdelbeginFL}{} %DIF PREAMBLE
\providecommand{\DIFdelendFL}{} %DIF PREAMBLE
%DIF HYPERREF PREAMBLE %DIF PREAMBLE
\providecommand{\DIFadd}[1]{\texorpdfstring{\DIFaddtex{#1}}{#1}} %DIF PREAMBLE
\providecommand{\DIFdel}[1]{\texorpdfstring{\DIFdeltex{#1}}{}} %DIF PREAMBLE
\providecommand{\DIFscaledelfig}{0.5}
%DIF HIGHLIGHTGRAPHICS PREAMBLE %DIF PREAMBLE
\RequirePackage{settobox} %DIF PREAMBLE
\RequirePackage{letltxmacro} %DIF PREAMBLE
\newsavebox{\DIFdelgraphicsbox} %DIF PREAMBLE
\newlength{\DIFdelgraphicswidth} %DIF PREAMBLE
\newlength{\DIFdelgraphicsheight} %DIF PREAMBLE
% store original definition of \includegraphics %DIF PREAMBLE
\LetLtxMacro{\DIFOincludegraphics}{\includegraphics} %DIF PREAMBLE
\providecommand{\DIFaddincludegraphics}[2][]{{\color{blue}\fbox{\DIFOincludegraphics[#1]{#2}}}} %DIF PREAMBLE
\providecommand{\DIFdelincludegraphics}[2][]{% %DIF PREAMBLE
\sbox{\DIFdelgraphicsbox}{\DIFOincludegraphics[#1]{#2}}% %DIF PREAMBLE
\settoboxwidth{\DIFdelgraphicswidth}{\DIFdelgraphicsbox} %DIF PREAMBLE
\settoboxtotalheight{\DIFdelgraphicsheight}{\DIFdelgraphicsbox} %DIF PREAMBLE
\scalebox{\DIFscaledelfig}{% %DIF PREAMBLE
\parbox[b]{\DIFdelgraphicswidth}{\usebox{\DIFdelgraphicsbox}\\[-\baselineskip] \rule{\DIFdelgraphicswidth}{0em}}\llap{\resizebox{\DIFdelgraphicswidth}{\DIFdelgraphicsheight}{% %DIF PREAMBLE
\setlength{\unitlength}{\DIFdelgraphicswidth}% %DIF PREAMBLE
\begin{picture}(1,1)% %DIF PREAMBLE
\thicklines\linethickness{2pt} %DIF PREAMBLE
{\color[rgb]{1,0,0}\put(0,0){\framebox(1,1){}}}% %DIF PREAMBLE
{\color[rgb]{1,0,0}\put(0,0){\line( 1,1){1}}}% %DIF PREAMBLE
{\color[rgb]{1,0,0}\put(0,1){\line(1,-1){1}}}% %DIF PREAMBLE
\end{picture}% %DIF PREAMBLE
}\hspace*{3pt}}} %DIF PREAMBLE
} %DIF PREAMBLE
\LetLtxMacro{\DIFOaddbegin}{\DIFaddbegin} %DIF PREAMBLE
\LetLtxMacro{\DIFOaddend}{\DIFaddend} %DIF PREAMBLE
\LetLtxMacro{\DIFOdelbegin}{\DIFdelbegin} %DIF PREAMBLE
\LetLtxMacro{\DIFOdelend}{\DIFdelend} %DIF PREAMBLE
\DeclareRobustCommand{\DIFaddbegin}{\DIFOaddbegin \let\includegraphics\DIFaddincludegraphics} %DIF PREAMBLE
\DeclareRobustCommand{\DIFaddend}{\DIFOaddend \let\includegraphics\DIFOincludegraphics} %DIF PREAMBLE
\DeclareRobustCommand{\DIFdelbegin}{\DIFOdelbegin \let\includegraphics\DIFdelincludegraphics} %DIF PREAMBLE
\DeclareRobustCommand{\DIFdelend}{\DIFOaddend \let\includegraphics\DIFOincludegraphics} %DIF PREAMBLE
\LetLtxMacro{\DIFOaddbeginFL}{\DIFaddbeginFL} %DIF PREAMBLE
\LetLtxMacro{\DIFOaddendFL}{\DIFaddendFL} %DIF PREAMBLE
\LetLtxMacro{\DIFOdelbeginFL}{\DIFdelbeginFL} %DIF PREAMBLE
\LetLtxMacro{\DIFOdelendFL}{\DIFdelendFL} %DIF PREAMBLE
\DeclareRobustCommand{\DIFaddbeginFL}{\DIFOaddbeginFL \let\includegraphics\DIFaddincludegraphics} %DIF PREAMBLE
\DeclareRobustCommand{\DIFaddendFL}{\DIFOaddendFL \let\includegraphics\DIFOincludegraphics} %DIF PREAMBLE
\DeclareRobustCommand{\DIFdelbeginFL}{\DIFOdelbeginFL \let\includegraphics\DIFdelincludegraphics} %DIF PREAMBLE
\DeclareRobustCommand{\DIFdelendFL}{\DIFOaddendFL \let\includegraphics\DIFOincludegraphics} %DIF PREAMBLE
%DIF AMSMATHULEM PREAMBLE %DIF PREAMBLE
\makeatletter %DIF PREAMBLE
\let\sout@orig\sout %DIF PREAMBLE
\renewcommand{\sout}[1]{\ifmmode\text{\sout@orig{\ensuremath{#1}}}\else\sout@orig{#1}\fi} %DIF PREAMBLE
\makeatother %DIF PREAMBLE
%DIF COLORLISTINGS PREAMBLE %DIF PREAMBLE
\RequirePackage{listings} %DIF PREAMBLE
\RequirePackage{color} %DIF PREAMBLE
\lstdefinelanguage{DIFcode}{ %DIF PREAMBLE
%DIF DIFCODE_UNDERLINE %DIF PREAMBLE
  moredelim=[il][\color{red}\sout]{\%DIF\ <\ }, %DIF PREAMBLE
  moredelim=[il][\color{blue}\uwave]{\%DIF\ >\ } %DIF PREAMBLE
} %DIF PREAMBLE
\lstdefinestyle{DIFverbatimstyle}{ %DIF PREAMBLE
	language=DIFcode, %DIF PREAMBLE
	basicstyle=\ttfamily, %DIF PREAMBLE
	columns=fullflexible, %DIF PREAMBLE
	keepspaces=true %DIF PREAMBLE
} %DIF PREAMBLE
\lstnewenvironment{DIFverbatim}[1][]{\lstset{style=DIFverbatimstyle}}{} %DIF PREAMBLE
\lstnewenvironment{DIFverbatim*}[1][]{\lstset{style=DIFverbatimstyle,showspaces=true}}{} %DIF PREAMBLE
\lstset{extendedchars=\true,inputencoding=utf8}

%DIF END PREAMBLE EXTENSION ADDED BY LATEXDIFF

\begin{document}
\maketitle

\setstretch{2}
Juda Kaleta

Institute of History, Faculty of Arts, Charles University, Czech Republic

Email: juda.kaleta@ff.cuni.cz | ORCID: 0000-0003-0043-4039

\hypertarget{abstract}{%
\subsection{Abstract}\label{abstract}}

This study reports \DIFdelbegin \DIFdel{the first }\DIFdelend \DIFaddbegin \DIFadd{a }\DIFaddend Czech adaptation of the Hartmann and
Hasselhorn Historical Perspective Taking (HPT) instrument and examines
whether its politically charged scenario introduces
construct-irrelevant variance (CIV) related to students' ideological
attitudes. Among 293 Czech secondary students, confirmatory factor
analysis \DIFdelbegin \DIFdel{(CFA) confirmed the }\DIFdelend \DIFaddbegin \DIFadd{favored a correlated }\DIFaddend three-factor \DIFdelbegin \DIFdel{structure. Item Response
Theory (IRT)-based differential item functioning (DIF) }\DIFdelend \DIFaddbegin \DIFadd{representation. An item-level
bias }\DIFaddend analysis detected no \DIFdelbegin \DIFdel{biased items }\DIFdelend \DIFaddbegin \DIFadd{differential item functioning }\DIFaddend across ideology
groups \DIFaddbegin \DIFadd{with this sample and procedure}\DIFaddend . Multi-group \DIFdelbegin \DIFdel{CFA
demonstrated full scalar measurement invariance}\DIFdelend \DIFaddbegin \DIFadd{factor analysis
produced fit-index changes compatible with scalar invariance, but
small extreme groups and unmodelled clustering limit the strength of that
evidence}\DIFaddend . Historical knowledge
\DIFdelbegin \DIFdel{predicted HPT scores; ideology was unrelated}\DIFdelend \DIFaddbegin \DIFadd{correlated with the primary six-item contextualization composite (}\emph{\DIFadd{r}} \DIFadd{= .34, 95\% CI }[\DIFadd{.23, .43}]\DIFadd{); the combined ideology score was
unrelated (}\emph{\DIFadd{r}} \DIFadd{= $-$.05, 95\% CI }[\DIFadd{$-$.16, .07}]\DIFadd{)}\DIFaddend . These findings support a
\DIFdelbegin \DIFdel{transferable }\DIFdelend \DIFaddbegin \DIFadd{possible }\DIFaddend protocol for evaluating measurement fairness when assessment
content carries attitudinal valence\DIFaddbegin \DIFadd{, subject to the limits of this
single scenario, item format, and ideologically restricted sample}\DIFaddend .

\textbf{Keywords:} historical perspective taking, cross-cultural validation, measurement invariance, differential item functioning, construct-irrelevant variance, test fairness

\hypertarget{introduction}{%
\subsection{1. Introduction}\label{introduction}}

Cross-cultural adaptation of educational and psychological instruments
requires more than linguistic translation: it demands evidence that
the adapted version measures the same construct, with equivalent
psychometric properties, in the target population \citep{AERA2014Standards,Hambleton2005,ITC2017}. When assessment content carries
attitudinal valence---when scenarios, stimuli, or items touch on
beliefs that respondents may hold---an additional fairness question
arises: does the content introduce construct-irrelevant variance (CIV)
by functioning differently for subgroups whose attitudes align with
that content \citep{Messick1995Validity}? This question is especially pressing for
scenario-based instruments in the social sciences, where rich
contextual material is needed to elicit complex reasoning but may
simultaneously activate respondents' pre-existing attitudes.
Historical perspective taking (HPT)---the capacity to reconstruct the
beliefs, values, and circumstances that made past actors' decisions
intelligible in their own time \DIFdelbegin %DIFDELCMD < \citep{Hartmann2008,Lee2001Empathy}%%%
\DIFdelend \DIFaddbegin \citep{Hartmann_Hasselhorn2008Historical,Lee_Ashby2001Empathy}\DIFaddend ---provides a particularly informative case. HPT
instruments routinely use politically and morally charged historical
scenarios as the stimulus for disciplinary reasoning (reasoning
according to the methodological norms of the historical discipline,
such as source criticism and \DIFdelbegin \DIFdel{contextualisation}\DIFdelend \DIFaddbegin \DIFadd{contextualization}\DIFaddend ; \citealp{VanDrie_VanBoxtel2008Historical}; \citealp{Seixas_Morton2013TheBig}; \citealp{Wineburg2001Historical}).

The instrument \DIFdelbegin \DIFdel{developed by }%DIFDELCMD < \citet{Hartmann2008} %%%
\DIFdel{is among
the most extensively validated standardized HPTmeasures in European
research. Originally }\DIFdelend \DIFaddbegin \DIFadd{reported by }\citet{Hartmann_Hasselhorn2008Historical} \DIFadd{provides a
scenario-based measure of HPT. It was }\DIFaddend validated with German \DIFdelbegin \DIFdel{secondary students , it has
been adapted for use in the Netherlands }%DIFDELCMD < \citep{Huijgen_et_al2014Testing,Huijgen_et_al2017Toward} %%%
\DIFdel{and employed in comparative studies of historical
reasoning}\DIFdelend \DIFaddbegin \DIFadd{Grade 10
students }\citep{Hartmann_Hasselhorn2008Historical} \DIFadd{and subsequently
translated and studied with Dutch students aged 10--17
}\citep{Huijgen_et_al2014Testing,Huijgen_et_al2017Toward}\DIFadd{. This evidence
base is promising but remains concentrated in two language contexts
and a small number of studies}\DIFaddend . The instrument operationalizes HPT through a scenario-based
approach in which students read about a historically situated dilemma
and respond to nine Likert-scaled items \DIFdelbegin \DIFdel{capturing three developmental
levels: presentism/othering of people }\DIFdelend \DIFaddbegin \DIFadd{in three interpretive
response categories: present-oriented perspective }\DIFaddend (POP), \DIFdelbegin \DIFdel{recognition of the
agent's social role and situational context }\DIFdelend \DIFaddbegin \DIFadd{role of the
historical agent }\DIFaddend (ROA), and \DIFaddbegin \DIFadd{historical }\DIFaddend contextualization (CONT). \DIFdelbegin \DIFdel{This structure reflects a theoretical progression from naive }\DIFdelend \DIFaddbegin \DIFadd{POP
and CONT represent the clearest contrast between }\DIFaddend present-centered
judgment \DIFdelbegin \DIFdel{toward increasingly sophisticated historical reasoning}%DIFDELCMD < \citep{Hartmann2008,Lee2001Empathy}%%%
\DIFdelend \DIFaddbegin \DIFadd{and contextual reasoning; ROA's position as an intermediate
level remains theoretically less certain
}\citep{Hartmann_Hasselhorn2008Historical,Lee_Ashby2001Empathy}\DIFaddend .

The present study reports \DIFdelbegin \DIFdel{the first }\DIFdelend \DIFaddbegin \DIFadd{a }\DIFaddend Czech adaptation and psychometric evaluation of the Hartmann and Hasselhorn HPT instrument. Two considerations motivate this work. First, no validated Czech-language HPT instrument \DIFdelbegin \DIFdel{currently exists, despite growing interest in historical thinking assessment within Czech education reform }%DIFDELCMD < \citep{Cinatl_et_al2021Historicka,Munich2023Zaverecna,Ripka_et_al2024Epistemic}%%%
\DIFdel{. The Czech adaptation addresses this gap and provides the research community with a tool for cross-cultural comparisons of HPT development}\DIFdelend \DIFaddbegin \DIFadd{was identified. Czech initiatives have developed digital historical-literacy tasks and measures of epistemic beliefs about history }\citep{Cinatl_et_al2021Historicka,Ripka_et_al2024Epistemic}\DIFadd{; the present adaptation provides a candidate instrument for further Czech validation and future cross-national equivalence studies}\DIFaddend .

Second, the instrument's scenario content raises a specific fairness question. \DIFdelbegin \DIFdel{The scenario describes a young man in the Weimar Republic deliberating whether to join an extremist political movement. Items that require students to contextualize this figure's reasoning }\DIFdelend \DIFaddbegin \DIFadd{Hannes discusses Germany's crisis and the 1930 election; students judge how well nine explanations fit his situation, including whether he could vote for an anti-democratic party such as the NSDAP. These judgments }\DIFaddend could, in principle, function differently for students whose political attitudes align with the scenario's ideological content. Students with right-authoritarian sympathies might find it cognitively easier to endorse contextualizing responses---not because they are reasoning historically, but because the responses resonate with their own attitudes. If present, such a pattern would constitute construct-irrelevant variance (CIV): systematic score variation attributable to a factor extraneous to the target construct \citep{AERA2014Standards,Messick1995Validity}.

This concern is not merely theoretical. Research on identity-protective cognition demonstrates that individuals process information through the lens of pre-existing beliefs \citep{Jost2003,Kahan2017}. Recent work has documented how ideological content can infiltrate competence judgments in history assessment, whether through teacher scoring bias in open-ended narrative
assessment \citep{Alven2019Bias} or through construct comparability challenges in cross-cultural settings \citep{Badham_et_al2025Construct}. The question central to the present study is whether the HPT instrument's psychometric properties---its factor structure, item parameters, and score interpretations---hold equivalently across students with different political attitudes, or whether politically charged content introduces measurement bias.

The study contributes to the literature in three ways: (a) it \DIFdelbegin \DIFdel{demonstrates a transferable }\DIFdelend \DIFaddbegin \DIFadd{applies a possible }\DIFaddend methodological approach for evaluating whether attitudinally charged assessment content introduces construct-irrelevant variance, (b) it provides DIF and \DIFdelbegin \DIFdel{measurement invariance evidence for a fairness-relevant grouping variable---political ideology---that has not previously been examined in HPT measurement}\DIFdelend \DIFaddbegin \DIFadd{measurement-invariance evidence across political-ideology groups}\DIFaddend , and (c) it \DIFdelbegin \DIFdel{reports the first published Czech validation of the Hartmann and Hasselhorn HPT instrument, extending the }\DIFdelend \DIFaddbegin \DIFadd{adds Czech evidence to the small }\DIFaddend cross-cultural evidence base for this measure. To establish this evidence, we conducted confirmatory factor analysis (CFA), reliability analysis, intraclass correlation estimation, IRT-based differential item functioning (DIF) analysis, and multi-group CFA for measurement invariance testing.

\hypertarget{literature-review}{%
\subsection{2. Literature Review}\label{literature-review}}

\hypertarget{historical-perspective-taking-and-its-measurement}{%
\subsubsection{2.1 Historical Perspective Taking and Its Measurement}\label{historical-perspective-taking-and-its-measurement}}

HPT refers to the ability to understand past actors' decisions by
reconstructing the historical context in which they operated \DIFdelbegin %DIFDELCMD < \citep{Hartmann2008,Lee2001Empathy}%%%
\DIFdelend \DIFaddbegin \citep{Hartmann_Hasselhorn2008Historical,Lee_Ashby2001Empathy}\DIFaddend . HPT is a second-order
(procedural or disciplinary) concept---as distinct from first-order
(substantive) concepts such as ``revolution'' or ``empire,'' it
concerns the meta-level reasoning processes historians use to make
sense of the past \DIFdelbegin %DIFDELCMD < \citep{Seixas2004Teaching}%%%
\DIFdelend \DIFaddbegin \citep{Seixas_Peck2004Teaching}\DIFaddend . The concept has roots in
British traditions that progressively refined the construct from
historical empathy \citep{Ashby1987,Shemilt1984} toward
rational reconstruction \DIFdelbegin %DIFDELCMD < \citep{Lee2001Empathy}%%%
\DIFdelend \DIFaddbegin \citep{Lee_Ashby2001Empathy}\DIFaddend . It has been situated
within broader frameworks of historical reasoning that encompass
evidence use, argumentation, and substantive concept deployment---the
application of domain-specific content knowledge about events, causes,
and actors \citep{vanBoxtel_vanDrie2018Historical,VanDrie_VanBoxtel2008Historical,VanSledright2014Assessing}. The literature variously labels this
capacity ``historical empathy,'' ``perspective taking,'' or ``rational
reconstruction.'' We follow \DIFdelbegin %DIFDELCMD < \citet{Hartmann2008} %%%
\DIFdelend \DIFaddbegin \citet{Hartmann_Hasselhorn2008Historical} \DIFaddend in adopting
the term ``perspective taking'' within the rationalist tradition of
\DIFdelbegin %DIFDELCMD < \citet{Lee2001Empathy}%%%
\DIFdel{, which emphasises }\DIFdelend \DIFaddbegin \citet{Lee_Ashby2001Empathy}\DIFadd{, which emphasizes }\DIFaddend cognitive reconstruction of
context rather than affective identification with past actors (see
\citealp{Barton2004}, for a contrasting view that emphasises
affective engagement alongside cognitive reconstruction; see also
\DIFdelbegin %DIFDELCMD < \citealp{Endacott2014Negotiating}%%%
\DIFdel{, for
a more recent treatment }\DIFdelend \DIFaddbegin \citealp{Endacott_Brooks2013An_Updated,Endacott2014Negotiating}\DIFadd{, for
treatments }\DIFaddend of the empathy/perspective-taking distinction). \DIFaddbegin \DIFadd{The present
instrument therefore operationalises cognitive perspective taking, not
the fuller cognitive-affective construct of historical empathy. }\DIFaddend Across these frameworks, HPT
is understood as a cognitively demanding, knowledge-dependent
competence---not a disposition but a disciplinary achievement requiring
both domain-specific knowledge and the capacity to deploy that
knowledge in reconstructing past perspectives \citep{Chapman2021Introduction,VanSledright2014Assessing,Wineburg2001Historical}. Distinguishing genuine
disciplinary reasoning from attitude-driven responding is a
measurement challenge that extends beyond HPT to epistemic cognition in
history more broadly \DIFdelbegin %DIFDELCMD < \citep{Maggioni2009Walking}%%%
\DIFdelend \DIFaddbegin \citep{Maggioni_et_al2009Walking}\DIFaddend .

\DIFdelbegin %DIFDELCMD < \citet{Hartmann2008} %%%
\DIFdel{operationalised HPT into a developmental
model with three levels: }\DIFdelend \DIFaddbegin \citet{Hartmann_Hasselhorn2008Historical} \DIFadd{operationalized HPT using three response
categories: }\DIFaddend \emph{\DIFdelbegin \DIFdel{presentism}\DIFdelend \DIFaddbegin \DIFadd{present-oriented perspective}\DIFaddend } (POP), where students judge
historical actors by contemporary standards; \emph{\DIFdelbegin \DIFdel{recognition }\DIFdelend \DIFaddbegin \DIFadd{role }\DIFaddend of the
\DIFaddbegin \DIFadd{historical }\DIFaddend agent\DIFdelbegin \DIFdel{'s role}\DIFdelend } (ROA), where students \DIFdelbegin \DIFdel{acknowledge the agent's social position without full contextual reconstruction}\DIFdelend \DIFaddbegin \DIFadd{explain action through familiar
or stereotyped social roles}\DIFaddend ; and
\emph{contextualization} (CONT), where students situate decisions
within the specific historical conditions that rendered them
intelligible. \DIFdelbegin \DIFdel{This developmental progression from present-centred
judgment
toward contextualisation }\DIFdelend \DIFaddbegin \DIFadd{The POP--CONT continuum from present-centered judgment
toward contextualization }\DIFaddend has empirical parallels in earlier
work on students' historical understanding \DIFdelbegin %DIFDELCMD < \citep{Lee2003A_scaffold}%%%
\DIFdelend \DIFaddbegin \citep{Lee_Shemilt2003A_scaffold}\DIFaddend . The
original \DIFaddbegin \DIFadd{published }\DIFaddend validation study (\emph{n} = 170 German students,
Grade 10\DIFdelbegin \DIFdel{)
recovered three levels through principal component
analysis }\DIFdelend \DIFaddbegin \DIFadd{; }\citealp{Hartmann_Hasselhorn2008Historical}\DIFadd{)
placed raw POP and CONT at opposing poles of one principal component
and ROA on a second, less clearly interpreted component }\DIFaddend (explained
variance = 51\%)\DIFdelbegin \DIFdel{, and students classified at the highest level }\DIFdelend \DIFaddbegin \DIFadd{; a separate latent-class analysis identified three
student classes. Students classified in the highest class }\DIFaddend had
significantly better history grades; however, PCA is an exploratory
technique, and the present study \DIFdelbegin \DIFdel{provides the first confirmatory factor
analytic test of the three-level structure.
The instrument was extended
by }\DIFdelend \DIFaddbegin \DIFadd{tests a correlated three-factor
representation rather than reproducing the original PCA dimensions.
}\DIFaddend \citet{Huijgen_et_al2014Testing} \DIFdelbegin \DIFdel{and }\DIFdelend \DIFaddbegin \DIFadd{replicated the two-dimensional pattern
for the Nazi-party scenario with Dutch students, but transfer of the
item-rating approach to a slavery scenario produced weaker reliability
and validity evidence. }\DIFaddend \citet{Huijgen_et_al2017Toward} \DIFdelbegin \DIFdel{, who confirmed the core theoretical structure while highlighting the difficulty of contextualization for most students}\DIFdelend \DIFaddbegin \DIFadd{subsequently
combined the closed item ratings with think-alouds and documented how
students interpreted and sometimes misunderstood the response options.
Thus, successful reuse of the Nazi scenario does not establish that the
format transfers unchanged across topics or elicits the intended process}\DIFaddend .

The difficulty of measuring HPT stems from its inherently interpretive character. Unlike factual recall, HPT cannot be \DIFdelbegin \DIFdel{assessed through simple multiple-choice items with single correct answers }\DIFdelend \DIFaddbegin \DIFadd{adequately inferred from decontextualized factual-recall items alone }\DIFaddend \citep{Breakstone2014Try,Smith2017Cognitive}. Scenario-based instruments address this challenge by presenting contextual information and requiring students to evaluate statements reflecting different levels of historical reasoning \DIFdelbegin %DIFDELCMD < \citep{Hartmann2008,Seixas_et_al2015A_Design}%%%
\DIFdelend \DIFaddbegin \citep{Hartmann_Hasselhorn2008Historical,Seixas_et_al2015A_Design}\DIFaddend . However, this approach introduces a measurement tension: the richer the historical context, the more opportunity for construct-irrelevant factors---including attitudes toward that content---to influence responses.

\hypertarget{cross-cultural-test-adaptation}{%
\subsubsection{2.2 Cross-Cultural Test Adaptation}\label{cross-cultural-test-adaptation}}

The adaptation of educational and psychological instruments across languages and cultures requires systematic procedures to ensure that the adapted version measures the same construct as the original \citep{Hambleton2005,ITC2017}. The International Test Commission (ITC) guidelines emphasize that adaptation involves more than translation: it requires attention to construct equivalence, cultural appropriateness, and technical quality of the adapted version \citep{ITC2017}. For instruments used in history education, cross-cultural adaptation faces the additional challenge that historical content may carry different resonances across national contexts. A scenario about the Weimar Republic, originally designed for German
students, intersects differently with Czech historical memory---while
subsequent events shape contemporary Czech memory of the interwar
period, the scenario centres on Weimar-era deliberation (early 1930s),
which predates the 1938 Munich Agreement and cession of the
Sudetenland by approximately \DIFdelbegin \DIFdel{five }\DIFdelend \DIFaddbegin \DIFadd{eight }\DIFaddend years. Nevertheless, the broader
period connects directly to the Nazi occupation of 1939--1945 in Czech
collective memory.

\DIFaddbegin \DIFadd{In Czech lower-secondary education, history is prescribed within the
compulsory People and Society curriculum area for Grades 6--9. The
framework includes the 1920s and 1930s, fascism, Nazism, totalitarianism,
and work with historical sources, but does not define HPT as a separate
competence }\citep[pp.~54--58, 140, 143]{RVPZV2023}\DIFadd{. School curricula broadly align
with this framework, although curriculum documents describe intended
rather than enacted provision }\citep{Kaleta2025Sticking}\DIFadd{. Participating
teachers reported that students had encountered the relevant interwar
content; this does not establish equivalent prior instruction across
schools.
}

\DIFadd{This documentary account applies specifically to the lower-secondary
subsample. The majority of participants were enrolled in upper-secondary
education, for which this study did not conduct an equivalent curriculum-
document audit. Eligibility at both levels required teacher confirmation that
students had completed the interwar unit, but that confirmation is not evidence
of uniform curricular coverage or instruction.
}

\DIFaddend Cross-cultural validation studies should demonstrate that the adapted instrument preserves the factor structure, reliability, and validity relationships of the original \citep{Hambleton2005}. At minimum, this requires confirming the theorized factor structure through CFA, reporting reliability estimates appropriate for the item characteristics, and providing convergent validity evidence through expected correlations with related constructs. When the instrument will be used to compare groups or predict outcomes, measurement invariance testing becomes essential \citep{AERA2014Standards}.

\hypertarget{construct-irrelevant-variance-in-educational-assessment}{%
\subsubsection{2.3 Construct-Irrelevant Variance in Educational Assessment}\label{construct-irrelevant-variance-in-educational-assessment}}

The unified validity framework identifies construct-irrelevant variance as a principal threat to score interpretation \citep{AERA2014Standards,Messick1995Validity}. CIV occurs when systematic variance in test scores is attributable to factors extraneous to the target construct. When assessment content intersects with respondents' attitudes or social identities, CIV can arise through attitude-consistent responding: students whose beliefs align with test content may find certain response options more natural or appealing, inflating their scores without corresponding gains in the assessed competence \citep{Messick1995Validity}.

For the HPT instrument specifically, the concern is that the Weimar Republic scenario creates a \emph{congruence pathway}---a mechanism whereby attitudinal similarity to item content lowers the cognitive effort required to endorse a response, inflating scores independently of competence. Consider a contextualization item such as ``His decision was understandable given the political and economic turmoil of the Weimar Republic.'' A student who endorses this statement because she has reconstructed the historical context produces a valid HPT response. But a student who endorses it because the sentiment resonates with her own political sympathies---not because she has engaged in historical reasoning---produces an inflated score. This would represent construct contamination in \citeauthor{Messick1995Validity}'s (\citeyear{Messick1995Validity}) terms---a systematic source of score variance unrelated to HPT competence. Evidence of internal structure, including measurement invariance across relevant subgroups, constitutes essential validity evidence when scores are used to compare groups or when the construct definition requires demonstrating that extraneous factors do not influence scores \citep{AERA2014Standards}.

An alternative theoretical prediction deserves mention. Research on authoritarian cognition suggests that high authoritarianism is associated with reduced integrative complexity---the cognitive capacity to hold and reconcile multiple conflicting perspectives simultaneously \citep{Duckitt2001,Suedfeld1977}. Because HPT requires precisely this kind of integrative capacity, authoritarian students might perform \emph{worse}, not better. The direction of any bias, if present, was genuinely uncertain. Additionally, dual-process theories of cognition \citep{Evans2013} suggest that structured assessment formats may channel deliberative processing, overriding the automatic attitude-retrieval processes that would produce attitude-consistent responding. If so, the assessment format itself may buffer against attitudinal interference.

\DIFaddbegin \DIFadd{The construct boundary is therefore consequential. Under broader accounts of
historical empathy, affective engagement and morally responsive evaluation can
be target-relevant aspects of understanding a past actor
}\citep{Barton2004,Endacott_Brooks2013An_Updated}\DIFadd{. Under the narrower cognitive
definition operationalized here, target variance consists of reconstructing
context and integrating historically relevant considerations. Moral reaction
or affect is not itself scored, and endorsement driven by option familiarity,
task demands, or attitude congruence is construct-irrelevant when it substitutes
for that reconstruction. The rating data alone cannot determine which process
produced an endorsement.
}

\DIFaddend \hypertarget{dif-and-measurement-invariance-as-fairness-evidence}{%
\subsubsection{2.4 DIF and Measurement Invariance as Fairness Evidence}\label{dif-and-measurement-invariance-as-fairness-evidence}}

Differential item functioning and measurement invariance testing provide complementary approaches to evaluating measurement fairness \citep{Chen2007,Cheung2002,Finch2005}. DIF analysis examines whether individual items function differently across groups after controlling for the latent trait, identifying specific items that may be biased. Measurement invariance testing through multi-group CFA examines whether the overall measurement model holds across groups at three progressively restrictive levels: \emph{configural invariance} (the same factors exist in both groups), \emph{metric invariance} (items relate equally strongly to the factors in both groups), and \emph{scalar invariance} (the scale origin is identical, so mean scores can be directly compared across groups).

For ordinal items, the weighted least squares mean and variance
adjusted (WLSMV) estimator---a robust estimation method designed for
categorical data---is recommended for CFA \citep{Flora2004}.
Because ordinal items have discrete response categories rather than
continuous scales, threshold-based invariance testing (which models the
boundaries between response categories) replaces intercept-based
approaches. The conventional criteria for evaluating invariance steps use changes in fit indices rather than chi-square difference tests, which are sensitive to sample size: \(\Delta\)CFI \(\leq\) .01 and \(\Delta\)RMSEA \(\leq\) .015 indicate that the more constrained model does not fit meaningfully worse \citep{Chen2007,Cheung2002}.

In the context of HPT measurement, establishing scalar invariance across ideology groups would mean that the factor structure, loadings, and item thresholds are equivalent for students with different political attitudes---a necessary condition for concluding that observed score differences (or the lack thereof) reflect true construct differences rather than measurement artifacts.

\hypertarget{the-present-study}{%
\subsubsection{2.5 The Present Study}\label{the-present-study}}

This study addresses two interrelated measurement questions:

\begin{enumerate}
\def\labelenumi{\arabic{enumi}.}
\item
  \DIFdelbegin \DIFdel{Does }\DIFdelend \DIFaddbegin \DIFadd{What psychometric evidence does }\DIFaddend the Czech adaptation of the Hartmann and Hasselhorn HPT instrument \DIFdelbegin \DIFdel{demonstrate acceptable psychometric properties, including the theorized }\DIFdelend \DIFaddbegin \DIFadd{provide, including evidence concerning a correlated }\DIFaddend three-factor \DIFdelbegin \DIFdel{structure, adequate }\DIFdelend \DIFaddbegin \DIFadd{representation, score }\DIFaddend reliability, and \DIFdelbegin \DIFdel{convergent validity }\DIFdelend \DIFaddbegin \DIFadd{convergence }\DIFaddend with historical knowledge?
\item
  \DIFdelbegin \DIFdel{Does the instrument function fairly across students with different political attitudes, as evidenced by the absence of DIFand by scalar measurement invariance across ideology groups}\DIFdelend \DIFaddbegin \DIFadd{What evidence do DIF, measurement-invariance, and regression analyses provide about whether scores are affected by students' political attitudes}\DIFaddend ?
\end{enumerate}

\DIFdelbegin \DIFdel{We hypothesized that the three-factor structure would be confirmed and that the instrument would demonstrate measurement invariance across ideology groups.
As convergent validity evidence, we expected historical knowledge to correlate positively with HPT scores, consistent with the theoretical conceptualization of HPT as
a knowledge-dependent competence}\DIFdelend \DIFaddbegin \DIFadd{The study and analysis plan were preregistered. The present
measurement-focused organization, the six-item contextualization
composite as the primary outcome, and several analyses were adopted
after registration, and not all registered analyses were implemented.
Table S5 maps the registered hypotheses, omissions, and
post-registration additions; findings from the latter are treated as
exploratory or descriptive rather than as preregistered tests}\DIFaddend .

\hypertarget{method}{%
\subsection{3. Method}\label{method}}

\hypertarget{adaptation-process}{%
\subsubsection{3.1 Adaptation Process}\label{adaptation-process}}

The Czech HPT instrument was developed through a systematic translation and adaptation process following ITC guidelines \citep{ITC2017}. Two bilingual Czech-English experts conducted independent forward translations from the \DIFdelbegin \DIFdel{published English version of the instrument}\DIFdelend \DIFaddbegin \DIFadd{German instrument reported in }\citet{Hartmann2008}\DIFadd{; wording was cross-checked against the published English validation article }\citep{Hartmann_Hasselhorn2008Historical}\DIFaddend . An independent backward translation was produced by a third bilingual expert unfamiliar with the original. Five Czech history educators at the secondary level reviewed the adapted items for conceptual equivalence, age-appropriateness, and curricular alignment. No items required substantive modification; minor wording adjustments were made to two items for naturalness in Czech. Reviewers \DIFdelbegin \DIFdel{confirmed that the Weimar Republic scenario content was familiar }\DIFdelend \DIFaddbegin \DIFadd{judged the scenario content accessible }\DIFaddend to Czech secondary students \DIFdelbegin \DIFdel{through the national history curriculum , given the direct connections between the interwar period and subsequent Czech historical experience (the Munich Agreement, the cession of the Sudetenland, and the Nazi occupation)}\DIFdelend \DIFaddbegin \DIFadd{given the curriculum coverage described in Section 2.2. This was expert-review evidence, not direct evidence from student interviews}\DIFaddend .

The adaptation preserved the \DIFdelbegin \DIFdel{original instrument's structure: the }\DIFdelend four-point response format (1 = \emph{strongly disagree} to 4 = \emph{strongly agree}), the \DIFdelbegin \DIFdel{three-factor item structure }\DIFdelend \DIFaddbegin \DIFadd{three item categories }\DIFaddend (three items per \DIFdelbegin \DIFdel{subscale}\DIFdelend \DIFaddbegin \DIFadd{category}\DIFaddend ), and \DIFdelbegin \DIFdel{the }\DIFdelend reverse-scoring of POP items \DIFdelbegin \DIFdel{were retained unchanged from the validated German version}\DIFdelend \DIFaddbegin \DIFadd{from the German source}\DIFaddend . All deviations from the original were documented transparently in the OSF repository (https://doi.org/10.17605/OSF.IO/YNG37). The \DIFdelbegin \DIFdel{scenario narrative was }\DIFdelend \DIFaddbegin \DIFadd{author separately }\DIFaddend cross-checked \DIFaddbegin \DIFadd{the scenario }\DIFaddend against the Czech history curriculum\DIFdelbegin \DIFdel{to confirm that all contextual references were accessible to the target population}\DIFdelend \DIFaddbegin \DIFadd{, but no direct student response-process study was conducted}\DIFaddend .

No cognitive pilot study (e.g., think-aloud interviews) was conducted prior to main data collection, which constitutes a limitation of the adaptation process (see Section 5.6).

\hypertarget{participants}{%
\subsubsection{3.2 Participants}\label{participants}}

The sample comprised 293 students from 20 classrooms across 10
schools in the Czech Republic, spanning lower-secondary
(International Standard Classification of Education [ISCED] level 2,
approximately ages 11--15; \emph{n} = 87, 29.7\%) and upper-secondary
(ISCED 3, approximately ages 15--19; \emph{n} = 206, 70.3\%) levels. The inclusion criterion was completion of the curricular unit on the interwar period. Schools were drawn from five Czech regions (Praha: 42.7\%; Královéhradecký: 34.8\%; Zlínský: 11.3\%; Jihomoravský: 8.2\%; Ústecký: 3.1\%) and three funding categories (public: 90.8\%; private: 6.5\%; church-affiliated: 2.7\%). Gender composition was 192 male (65.5\%), 88 female (30.0\%), and 13 identifying as another gender (4.4\%). Students' most recent history grade averaged 1.71 (\emph{SD} = 0.92; Czech scale: 1 = excellent, 5 = failing).

Schools were recruited through convenience sampling; the sample over-represents upper-secondary students relative to the national age cohort composition and should not be treated as a probability sample of Czech adolescents. Sample size was determined by an a priori power analysis specified in the preregistered protocol (OSF: https://doi.org/10.17605/OSF.IO/YNG37), targeting 80\% power to detect standardized ideology effects of \(\beta\) \(\geq\) .15 in a two-level multilevel \DIFdelbegin \DIFdel{model}\DIFdelend \DIFaddbegin \DIFadd{regression}\DIFaddend ; the achieved \emph{N} = 285--293 met \DIFdelbegin \DIFdel{this target. }\DIFdelend \DIFaddbegin \DIFadd{that regression target. The calculation did not establish power for the post-registration DIF or MG-CFA analyses. }\DIFaddend Within each school, all \DIFaddbegin \DIFadd{participating }\DIFaddend students present on the day of administration completed the instrument battery during a regular class period (approximately 30--40 minutes). Teachers were present during administration but did not intervene in students' responses.

\DIFaddbegin \DIFadd{The sample's right-authoritarian attitude median was 2.33 on the
five-point scale, and the historical-knowledge median was 3 of 6; fuller
distributions appear in Section 4.4. No defensible student-, family-, or
school-level socioeconomic indicator was collected. Consequently, the sample's
socioeconomic composition and any associated selection bias cannot be
characterized.
}

\DIFaddend \hypertarget{instruments}{%
\subsubsection{3.3 Instruments}\label{instruments}}

\hypertarget{historical-perspective-taking-hpt}{%
\paragraph{3.3.1 Historical Perspective Taking (HPT)}\label{historical-perspective-taking-hpt}}

Students completed the Czech adaptation of the \DIFdelbegin %DIFDELCMD < \citet{Hartmann2008} %%%
\DIFdelend scenario-based HPT instrument \DIFaddbegin \DIFadd{from the German dissertation }\citep{Hartmann2008}\DIFadd{, with the published English account used as supporting validation evidence }\citep{Hartmann_Hasselhorn2008Historical}\DIFaddend . They read a narrative \DIFdelbegin \DIFdel{about a young man in the Weimar Republic considering joining an extremist political movement, then responded to nine items }\DIFdelend \DIFaddbegin \DIFadd{in which Hannes discusses Germany's crisis and the 1930 election, then judged nine explanations of his situation, including whether he could vote for an anti-democratic party such as the NSDAP, }\DIFaddend on a 4-point Likert scale. Three items each tapped \DIFdelbegin \DIFdel{presentism }\DIFdelend \DIFaddbegin \DIFadd{present-oriented perspective }\DIFaddend (POP1--POP3), \DIFdelbegin \DIFdel{recognition of the agent's role }\DIFdelend \DIFaddbegin \DIFadd{role of the historical agent }\DIFaddend (ROA1--ROA3), and \DIFaddbegin \DIFadd{historical }\DIFaddend contextualization (CONT1--CONT3). POP items capture present-centered \DIFdelbegin \DIFdel{moral }\DIFdelend judgments of the historical figure\DIFdelbegin \DIFdel{(e.g., evaluating his decision by contemporary standards)}\DIFdelend ; these are reverse-scored (5 $-$ POP, yielding a reversed range of 1--4) so that higher scores indicate \DIFdelbegin \DIFdel{more contextualized }\DIFdelend \DIFaddbegin \DIFadd{less present-oriented }\DIFaddend reasoning. ROA items \DIFdelbegin \DIFdel{reflect awareness of the agent's social position and circumstances without full contextual integration}\DIFdelend \DIFaddbegin \DIFadd{explain action through familiar or stereotyped social roles}\DIFaddend . CONT items \DIFdelbegin \DIFdel{represent the most sophisticated level, requiring }\DIFdelend \DIFaddbegin \DIFadd{require }\DIFaddend students to situate the historical figure's deliberation within the specific political, economic, and social conditions of the Weimar Republic. The primary \DIFdelbegin \DIFdel{composite was HPT\_CTX6, }\DIFdelend \DIFaddbegin \DIFadd{six-item contextualization composite was }\DIFaddend the mean of reversed \DIFdelbegin \DIFdel{POP and CONT items. HPT\_TOT9 (adding ROA) served }\DIFdelend \DIFaddbegin \DIFadd{present-oriented perspective and contextualization scores and therefore represents the continuum between these two response categories. The nine-item total score (adding the theoretically less certain role-of-the-agent category) served only }\DIFaddend as a secondary \DIFaddbegin \DIFadd{sensitivity }\DIFaddend outcome.

\hypertarget{right-authoritarian-ideology-fr-lf-mini}{%
\paragraph{3.3.2 Right-Authoritarian Ideology (FR-LF mini)}\label{right-authoritarian-ideology-fr-lf-mini}}

A six-item short form of the Fragebogen Rechtsextremer Einstellungen---Leipziger Form (Questionnaire on Right-Wing Extremist Attitudes---Leipzig Form; \citealp{Decker_et_al2013Rechtsextremismus}), comprising acceptance of right-wing dictatorship (RD; 3 items) and National Socialist sympathy/relativization (NS; 3 items) facets, rated on a 5-point scale. The NS facet is theoretically the most direct operationalization of potential content congruence with the Weimar scenario.

\hypertarget{authoritarianism-ksa-3}{%
\paragraph{3.3.3 Authoritarianism (KSA-3)}\label{authoritarianism-ksa-3}}

The nine-item Kurzskala Autoritarismus (Short Scale of
Authoritarianism; \citealp{Beierlein_et_al2014Die_Kurzskala}), measuring authoritarian
aggression, submission, and conventionalism on a 5-point scale. This instrument was used for a supplementary measurement invariance check using an alternative ideology operationalization.

\hypertarget{historical-knowledge-kn}{%
\paragraph{3.3.4 Historical Knowledge (KN)}\label{historical-knowledge-kn}}

Six multiple-choice items covering content relevant to the HPT scenario, scored dichotomously (sum score 0--6). This measure served as a convergent validity criterion, given the theoretical expectation that HPT is a knowledge-dependent competence.

\hypertarget{social-desirability-sdr-5}{%
\paragraph{3.3.5 Social Desirability (SDR-5)}\label{social-desirability-sdr-5}}

The five-item short form \citep{Hays_et_al1989A_Five}, rated on a 5-point scale with items SDR2--SDR4 reverse-coded. Included to evaluate whether HPT scores reflect impression management rather than genuine reasoning. \DIFaddbegin \DIFadd{The source studies for the FR-LF included German respondents aged 14 years and older, whereas KSA-3 was developed for German-speaking adults aged 18 years and older and SDR-5 with adult medical outpatients and older adults. None therefore constitutes direct Czech adolescent validation; suitability rests on expert review and the sample-specific evidence reported here and in Table S4.
}\DIFaddend

\hypertarget{analytic-approach}{%
\subsubsection{3.4 Analytic Approach}\label{analytic-approach}}

All analyses were conducted in R (version 4.4.2) using lavaan for CFA, mirt for IRT-based DIF analysis, psych for reliability estimation, and lme4/performance for intraclass correlations.

\textbf{Factor structure.} CFA compared one-factor (general HPT), two-factor (POP+CONT vs.~ROA), and three-factor (POP, ROA, CONT) models using the WLSMV estimator, appropriate for ordinal item responses \citep{Flora2004}. CFA fit indices summarize how well a hypothesized factor structure reproduces the observed correlations among items; multiple indices are reported because each captures a different aspect of misfit. Model evaluation used CFI (\(\geq\) .95 acceptable, \(\geq\) .97
good), the Tucker--Lewis Index (TLI; \(\geq\) .95), the Root Mean
Square Error of Approximation (RMSEA; \(\leq\) .08 acceptable,
\(\leq\) .05 good), and the Standardized Root Mean Square Residual
(SRMR; \(\leq\) .08).

\textbf{Reliability.} Internal consistency was assessed \DIFdelbegin \DIFdel{via }\DIFdelend \DIFaddbegin \DIFadd{using raw and
polychoric }\DIFaddend Cronbach's \(\alpha\)\DIFdelbegin \DIFdel{(raw and
based on polychoric correlations) and }\DIFdelend \DIFaddbegin \DIFadd{, }\DIFaddend McDonald's \(\omega\), \DIFdelbegin \DIFdel{which
does not assume tau-equivalence (the assumption that all itemsrelate
equally strongly to the latent factor) and is generally preferred for
congeneric measures (measures whose items may differ in how strongly
they relate to the factor; }%DIFDELCMD < \citealp{Flora2004}%%%
\DIFdel{)}\DIFdelend \DIFaddbegin \DIFadd{and KR-20 for
the dichotomous knowledge items. Reliability of the combined ideology
score was estimated with the Mosier formula. Full methods and estimates
are reported in Table S3}\DIFaddend .

\textbf{Intraclass correlations.} We computed intraclass correlation
coefficients (ICCs) at the school and classroom levels to \DIFdelbegin \DIFdel{characterise
}\DIFdelend \DIFaddbegin \DIFadd{characterize
}\DIFaddend the multilevel structure of the data and to evaluate whether
clustering warranted multilevel approaches.

\textbf{Clustered data handling.} The \DIFdelbegin \DIFdel{DIF and }\DIFdelend \DIFaddbegin \DIFadd{substantive regressions used
random intercepts for the nested data. The DIF, }\DIFaddend MG-CFA\DIFdelbegin \DIFdel{analyses treat
}\DIFdelend \DIFaddbegin \DIFadd{, and zero-order
correlation analyses treated }\DIFaddend observations as independent\DIFdelbegin \DIFdel{, which is standard practice in applied DIF
and measurement invariance research. To evaluate whether ignoring the
nested structure introduces meaningful bias, design effects (DEFF---a
multiplier indicating how much clustering inflates standard errors relative to a simple random sample) were computed post hoc. A design
effect of 1.0 means clusteringhas no impact on precision; values
approaching 2.0 indicate modest inflation of standard errors that can
usually be tolerated }%DIFDELCMD < \citep{Maas2005}%%%
\DIFdel{. The primary composite (HPT\_CTX6) had DEFF = 1.56, within the conventional \(\leq\) 2.0 threshold.
The secondary composite (HPT
\_TOT9) had DEFF = 2.16, marginally above this threshold; consequently, inferential claims are anchored primarily to HPT\_CTX6, and HPT\_TOT9 results should be interpreted with this caveat.
}\footnote{\DIFdel{DEFF = 1 + (\(\bar{n}_c\) $-$ 1) \(\times\) ICC, where \(\bar{n}_c\) is the average cluster size (approximately 14.7 students per school). The TYPE=COMPLEX option in lavaan was not applied because the DIF analysis in mirt does not support clustered standard errors, and applying different corrections to DIF and MG-CFA would make the two sets of results less directly comparable.}}
%DIFAUXCMD
\addtocounter{footnote}{-1}%DIFAUXCMD
\DIFdelend \DIFaddbegin \DIFadd{. Composite
score ICCs do not establish that item thresholds, loadings, or their
standard errors are unaffected by clustering; therefore, uncertainty
for these analyses may be understated. The zero-order correlation
intervals are reported as descriptive, unadjusted intervals rather than
cluster-robust inference.
}\DIFaddend

\DIFaddbegin \textbf{\DIFadd{Missing data and score construction.}} \DIFadd{Composite scores required
at least two answered items per three-item HPT
subscale, four of six FR-LF items, seven of nine KSA-3 items, and four
of five SDR-5 items; available responses were averaged. Missing
knowledge responses were scored as incorrect, following the original
analysis scripts. Reliability estimates used listwise-complete item
sets, while correlations used pairwise-complete composite scores. These
rules explain the analysis-specific sample sizes reported in the tables
and supplement.
}

\DIFaddend \textbf{Differential item functioning.} Low and High ideology groups were formed \DIFaddbegin \DIFadd{post-registration }\DIFaddend using tertile splits on a composite ideology score (mean of \emph{z}-scored FR-LF and KSA-3), with the middle tertile excluded \DIFdelbegin \DIFdel{to maximize contrast (}\DIFdelend \DIFaddbegin \DIFadd{(Low }\DIFaddend \emph{n} \DIFdelbegin \DIFdel{\(\approx\) }\DIFdelend \DIFaddbegin \DIFadd{= }\DIFaddend 96\DIFdelbegin \DIFdel{per group). This extreme-group design follows established practice in DIF research }%DIFDELCMD < \citep{Finch2005}%%%
\DIFdel{. A constrained }\DIFdelend \DIFaddbegin \DIFadd{; High }\emph{\DIFadd{n}} \DIFadd{= 95). A constrained unidimensional }\DIFaddend multi-group graded response model (GRM item
discrimination and \DIFdelbegin \DIFdel{threshold }\DIFdelend \DIFaddbegin \DIFadd{category-intercept }\DIFaddend parameters for each item by group are
reported in Table S1 in the OSF supplementary materials) using a
sequential item-release design tested each item for \DIFdelbegin \DIFdel{uniform DIF (items
are systematically easier or harder for one group at all ability
levels) and non-uniform DIF (the group difference reverses depending
on ability level) }\DIFdelend \DIFaddbegin \DIFadd{omnibus DIF by
jointly freeing its discrimination and category intercepts, }\DIFaddend with
Bonferroni correction (\(\alpha\) = .01). The
baseline model constrains all nine items to equality across groups;
each test then frees the discrimination and \DIFdelbegin \DIFdel{threshold parameters }\DIFdelend \DIFaddbegin \DIFadd{category intercepts }\DIFaddend for
the focal item, yielding a likelihood-ratio comparison against the
fully constrained baseline (an all-others-as-anchor strategy). This
approach assumes that the remaining eight items are DIF-free when
testing each focal item; iterative purification was not applied, which
means that if multiple \DIFaddbegin \DIFadd{anchor }\DIFaddend items exhibited small DIF\DIFdelbegin \DIFdel{in the same
direction, anchor contamination could mask detection }%DIFDELCMD < \citep{Woods2009}%%%
\DIFdel{.
}\DIFdelend \DIFaddbegin \DIFadd{, focal-item
detection could be affected.
Because the CFA favoured correlated dimensions, the unidimensional GRM
is a pragmatic sensitivity model rather than a fully aligned
measurement model; multidimensionality or local dependence may affect
the item tests.
}\DIFaddend

\textbf{Sensitivity of DIF detection.} With approximately 96 respondents per ideology group \DIFdelbegin \DIFdel{, the likelihood-ratio DIF testsat Bonferroni-corrected \(\alpha\) = .01 (9 tests) are powered primarily to detect moderate-to-large parameter differences (\(\Delta\beta\) \(\geq\) 0.5 logits on discrimination or threshold parameters). Weak DIF effects (\(\Delta\beta\) $<$ 0.5 logits) cannot be confidently excluded. This sensitivity threshold is consistent with the sample size requirements documented for IRT-based DIF methods }%DIFDELCMD < \citep{Finch2005}%%%
\DIFdelend \DIFaddbegin \DIFadd{and nine Bonferroni-adjusted tests, statistical power was limited. No design-specific simulation was conducted, so the minimum detectable DIF magnitude is unknown}\DIFaddend . The extreme-group design \DIFdelbegin \DIFdel{maximizes statistical contrast between ideology subgroups }\DIFdelend \DIFaddbegin \DIFadd{increases ideological contrast }\DIFaddend at the cost of \DIFdelbegin \DIFdel{reduced per-group sample size---a standard trade-off in DIF research when the grouping variable is continuous}\DIFdelend \DIFaddbegin \DIFadd{excluding the middle tertile and reducing group size}\DIFaddend .

\textbf{Multi-group CFA.} Stepwise MG-CFA tested configural, metric, and scalar invariance across ideology groups using the three-factor model with WLSMV estimation. Invariance was evaluated using \(\Delta\)CFI \(\leq\) .01 and \(\Delta\)RMSEA \(\leq\) .015 criteria \citep{Chen2007,Cheung2002}.

\DIFaddbegin \textbf{\DIFadd{Substantive models.}} \DIFadd{Random-intercept linear models predicted each
standardized HPT outcome from standardized FR-LF, KSA-3, historical-knowledge,
and SDR-5 scores entered together, with intercepts for school and classroom.
Thus, each reported ideology coefficient is mutually adjusted for the other
ideology measure, knowledge, social desirability, and the nested grouping
structure.
}

\DIFaddend \textbf{Convergent validity.} \DIFaddbegin \DIFadd{Descriptive }\DIFaddend Pearson correlations between HPT composites and historical knowledge provided convergent validity evidence\DIFaddbegin \DIFadd{; Fisher-transformed intervals did not adjust for clustering}\DIFaddend .

\textbf{Supplementary invariance check.} To verify that invariance findings were not an artifact of the composite ideology definition, MG-CFA was also conducted using NS-only grouping (the facet most directly content-congruent with the Weimar scenario).

\hypertarget{results}{%
\subsection{4. Results}\label{results}}

\DIFaddbegin \DIFadd{The Results follow the two measurement questions in sequence. Sections
4.1--4.5 assess factor structure, score precision, distributions, and
clustering; Sections 4.6--4.8 examine item and model-level fairness; Sections
4.9--4.10 then report the focal knowledge and ideology relationships.
}

\DIFaddend \hypertarget{factor-structure}{%
\subsubsection{4.1 Factor Structure}\label{factor-structure}}

CFA supported the three-factor model (POP, ROA, CONT) as best-fitting, outperforming the two-factor and one-factor alternatives (Table 1). The three-factor model showed excellent global fit: CFI = .999, TLI = .998, RMSEA = .011 (90\% CI [.000, .051]), SRMR = .048. The two-factor model combining POP and CONT items on a single factor showed acceptable but inferior fit (CFI = .956, RMSEA = .057), while the one-factor model fit poorly (CFI = .926, RMSEA = .073).

\textbf{Table 1.} \emph{CFA Fit Indices for HPT Factor Models (WLSMV Estimator)}

% Table shown without inline diff markup for longtable stability.
\begin{longtable}[]{@{}lrrrrrll@{}}
\toprule
Model & \(\chi^2\) & \emph{df} & CFI & TLI & RMSEA & 90\% CI & SRMR \\
\midrule
\endhead
\bottomrule
\endlastfoot
1-Factor & 66.68 & 27 & .926 & .902 & .073 & [.051, .095] & .075 \\
2-Factor (POP+CONT, ROA) & 49.48 & 26 & .956 & .940 & .057 & [.032, .081] & .066 \\
3-Factor (POP, ROA, CONT) & 24.78 & 24 & .999 & .998 & .011 & [.000, .051] & .048 \\
\end{longtable}

\emph{Note.} WLSMV = weighted least squares mean and variance adjusted. All items treated as ordered categorical. The 2-factor model groups POP and CONT items on a single factor with ROA on a separate factor. Chi-square values are mean-and-variance adjusted (WLSMV). Nested model comparisons require the DIFFTEST procedure rather than standard chi-square difference tests. All \emph{p} $<$ .001 except the three-factor model (\emph{p} = .418).

The near-perfect global fit of the three-factor model reflects the
small item set (nine items in three just-identified three-indicator
factors---each subscale has exactly as many free parameters as observed
statistics, leaving no room for misfit within subscales) and should be
interpreted alongside local fit diagnostics: the largest modification
index was 7.35 (POP1--ROA3 residual covariance; below the threshold
of 10), and standardized residual correlations were small (\(|r|\)
\(\leq\) .13). With only \DIFdelbegin \emph{\DIFdel{df}} %DIFAUXCMD
\DIFdel{= 24 for the three-factor model (9
items, 3 factors, 24 free parameters against 45 observed statistics),
there is minimal room for }\DIFdelend \DIFaddbegin \DIFadd{three indicators per factor, there is limited
opportunity to reveal within-factor }\DIFaddend misfit; the near-perfect CFI is therefore a
structural consequence of the low-\emph{df} specification rather than
evidence of an unusually well-fitting model. \DIFdelbegin \DIFdel{These results confirm
that the Czech adaptation preserves the theorized three-level
structure of the original German instrument}\DIFdelend \DIFaddbegin \DIFadd{The Czech data are
compatible with the correlated three-factor specification, but this is
not a test of cross-national equivalence}\DIFaddend .

Inter-factor correlations from the three-factor CFA were: POP--ROA \(\Phi\) = .377, POP--CONT \(\Phi\) = .509, ROA--CONT \(\Phi\) = .631. These moderate-to-strong correlations \DIFdelbegin \DIFdel{indicate meaningful associations among the HPT dimensions, with ROA and CONT showing the strongest link. Discriminant validityis supported because all inter-factor
correlations are below .85,
confirming that the three subscales capture related but distinct
aspects of historical perspective taking}\DIFdelend \DIFaddbegin \DIFadd{and the better fit of the three-factor model provide internal-structure evidence that the item sets are distinguishable in this sample. They do not by themselves establish discriminant validity. Stronger evidence would require showing that the subscales relate differentially to external criteria in theoretically expected ways}\DIFaddend .

A supplementary bifactor model (one general factor plus three specific
factors, in which each item loads on both the general factor and its
subscale-specific factor) fit well (\(\chi^2\) = 14.70, \emph{df} =
18, \emph{p} = .682) but did not meaningfully improve over the
correlated three-factor model (\(\Delta\)CFI = .001). The low
\emph{df} (18) results from the additional cross-loadings consuming
degrees of freedom, and the comparison with the correlated three-factor
model (\emph{df} = 24) should be treated cautiously given the
different parameterisations. CONT items loaded primarily on the
general factor (.47--.63), while POP and ROA retained meaningful
specific factor variance, supporting subscale-level interpretation.

\hypertarget{standardized-factor-loadings}{%
\subsubsection{4.2 Standardized Factor Loadings}\label{standardized-factor-loadings}}

Table 2 presents the standardized loadings from the three-factor CFA. CONT items loaded most strongly (.61--.72), followed by POP items (.30--.74) and ROA items (.30--.67). POP2 showed the weakest loading (.30), at the conventional minimum threshold for item retention; it was retained to preserve fidelity to the original instrument structure.
\DIFdelbegin \DIFdel{The range of loadings across subscales reflects the heterogeneity within each three-item set and is consistent with patterns reported in prior validations of short-form instruments.
}\DIFdelend

\textbf{Table 2.} \emph{\DIFaddbegin \DIFadd{Summary of }\DIFaddend Standardized Factor Loadings\DIFdelbegin \DIFdel{for the Three-Factor HPT Model}\DIFdelend }

% Table shown without inline diff markup for longtable stability.
\begin{longtable}[]{@{}lll@{}}
\toprule
Subscale & Loading range & Weakest item \\
\midrule
\endhead
\bottomrule
\endlastfoot
Present-oriented perspective & .30--.74 & POP2 (.30) \\
Role of historical agent & .30--.67 & ROA2 (.30) \\
Historical contextualization & .61--.72 & CONT2 (.61) \\
\end{longtable}

\emph{Note.} POP items are reverse-scored (higher = more contextualized reasoning). All loadings \DIFaddbegin \DIFadd{were }\DIFaddend significant at \emph{p} $<$ .001. \DIFdelbegin \DIFdel{WLSMV estimator with items treated as ordered categorical}\DIFdelend \DIFaddbegin \DIFadd{Full item estimates, standard errors, and confidence intervals are reported in Table S2}\DIFaddend .

\hypertarget{reliability-and-internal-consistency}{%
\subsubsection{4.3 Reliability and Internal Consistency}\label{reliability-and-internal-consistency}}

Table 3 \DIFdelbegin \DIFdel{presents reliability coefficients for each HPT subscale. Raw }\DIFdelend \DIFaddbegin \DIFadd{summarizes reliability for the primary study scores. Across HPT subscales, raw }\DIFaddend Cronbach's \(\alpha\) \DIFdelbegin \DIFdel{values ranged from .47 (POP) }\DIFdelend \DIFaddbegin \DIFadd{ranged from .46 }\DIFaddend to .64 \DIFdelbegin \DIFdel{(CONT); polychoric \(\alpha\) (appropriate for ordinal items) ranged from .53 to .69; }\DIFdelend \DIFaddbegin \DIFadd{and }\DIFaddend McDonald's \(\omega\) \DIFdelbegin \DIFdel{ranged from .55 to .69. CONT demonstrated the strongest reliability, consistent with its highest average factor loadings}\DIFdelend \DIFaddbegin \DIFadd{from .54 to .70. Full estimates for all HPT and predictor scores are in Table S3}\DIFaddend .

\textbf{Table 3.} \emph{Reliability \DIFdelbegin \DIFdel{Coefficients }\DIFdelend \DIFaddbegin \DIFadd{Summary }\DIFaddend for \DIFdelbegin \DIFdel{HPT Subscales}\DIFdelend \DIFaddbegin \DIFadd{Primary Study Scores}\DIFaddend }

% Table shown without inline diff markup for longtable stability.
\begin{longtable}[]{@{}lccc@{}}
\toprule
Score & Items/components & Reliability & Method \\
\midrule
\endhead
\bottomrule
\endlastfoot
HPT composite & 6 & .66 & \(\omega\) total \\
HPT total & 9 & .70 & \(\omega\) total \\
Combined ideology & 2 scales & .80 & Mosier \\
\end{longtable}

\emph{Note.} \DIFdelbegin \DIFdel{POP items are reverse-scored. }\DIFdelend \DIFaddbegin \DIFadd{The Mosier estimate has a 95\% bootstrap CI }[\DIFadd{.74, .84}]\DIFadd{. Full subscale and predictor alpha, polychoric-alpha, omega, and interval estimates are reported in Table S3.
}

\DIFadd{Reliability was modest and constrains individual interpretation. FR-LF
had }\DIFaddend \(\alpha\) = \DIFdelbegin \DIFdel{Cronbach's alpha; \(\omega\) }\DIFdelend \DIFaddbegin \DIFadd{.67, KSA-3 had \(\alpha\) }\DIFaddend = \DIFdelbegin \DIFdel{McDonald's omegatotal.
Mean \(r_{ii}\) }\DIFdelend \DIFaddbegin \DIFadd{.72, and the combined
ideology score had estimated reliability .80, 95\% CI }[\DIFadd{.74, .84}]\DIFadd{.
Knowledge had KR-20 }\DIFaddend = \DIFdelbegin \DIFdel{mean inter-item correlation.
}%DIFDELCMD <

%DIFDELCMD < %%%
\DIFdel{These values are modest by conventional standards but structurally constrained by the three-item design: with }\emph{\DIFdel{k}} %DIFAUXCMD
\DIFdel{= 3, even moderate inter-item correlations (.20--.35) yield }\DIFdelend \DIFaddbegin \DIFadd{.55, and SDR-5 was weak (}\DIFaddend \(\alpha\) \DIFdelbegin \DIFdel{in this range. The Spearman-Brown formula (which estimates how reliability would improve if subscale length were doubled) projects \(\alpha\) }\DIFdelend = \DIFdelbegin \DIFdel{.64--.79 for six-item subscales, though this assumes tau-equivalence---an assumption violated when loadings range from }\DIFdelend .30\DIFdelbegin \DIFdel{to .74 (Table 2}\DIFdelend )\DIFdelbegin \DIFdel{---so the projections should be treated as approximate upper bounds. McDonald's \(\omega\) (.55--.69) is more appropriate given unequal loadings (Table 2). The knowledge--HPT correlations reported below (\(r\) = .22--.34) demonstrate that the instrument carries genuine psychometric signal despite modest subscale reliability }\DIFdelend .
\DIFdelbegin \DIFdel{Correcting the ideology--HPT correlation for measurement unreliability
(disattenuation---an adjustment that estimates what the correlation
would be if both measures had perfect reliability) yielded a
disattenuated \(r\) = -.07, indicating that measurement error does not
mask a meaningful ideology association. However, reliability for the
composite ideology score (the mean of }\emph{\DIFdel{z}}%DIFAUXCMD
\DIFdel{-scored FR-LF and KSA-3) was not separately estimated, which limits the precision of
this disattenuated estimate.
}\DIFdelend

\hypertarget{descriptive-statistics}{%
\subsubsection{4.4 Descriptive Statistics}\label{descriptive-statistics}}

Table 4 presents descriptive statistics. HPT scores followed the expected \DIFdelbegin \DIFdel{developmental ordering(POP\_rev $>$ ROA $>$ CONT), with contextualization producing }\DIFdelend \DIFaddbegin \DIFadd{mean-score ordering: reversed present-oriented perspective scores were highest, followed by role-of-the-agent and contextualization scores. Contextualization produced }\DIFaddend the lowest scores (\emph{M} = 2.71)\DIFdelbegin \DIFdel{. Ideology scores were below the scale midpoint (}\DIFdelend \DIFaddbegin \DIFadd{; this ordering may reflect item endorsement or difficulty and is not developmental evidence. Right-authoritarian attitude scores (}\DIFaddend FR-LF\DIFdelbegin \DIFdel{: }\DIFdelend \DIFaddbegin \DIFadd{) and authoritarianism scores (KSA-3) were below their scale midpoints (}\DIFaddend \emph{M} = 2.48 \DIFdelbegin \DIFdel{; KSA: }\emph{\DIFdel{M}} %DIFAUXCMD
\DIFdel{= }\DIFdelend \DIFaddbegin \DIFadd{and }\DIFaddend 2.86\DIFdelbegin \DIFdel{), reflecting the rarity of strong }\DIFdelend \DIFaddbegin \DIFadd{, respectively). The }\DIFaddend right-authoritarian \DIFdelbegin \DIFdel{endorsement in mainstream adolescent samples. }\DIFdelend \DIFaddbegin \DIFadd{attitude median was 2.33 (IQR }[\DIFadd{2.00, 3.00}]\DIFadd{) and 51 of 284 students (18.0\%) scored above its midpoint. The authoritarianism median was 2.89 (IQR }[\DIFadd{2.53, 3.22}]\DIFadd{) and 106 of 283 students (37.5\%) scored above its midpoint. On the mean of the two raw scales, 71 of 283 students (25.1\%) scored above 3. Historical knowledge had a median of 3 (IQR }[\DIFadd{2, 4}]\DIFadd{), with 4.1\% at the floor and 8.9\% at the ceiling. No defensible student- or school-level socioeconomic measure was available.
}\DIFaddend

\textbf{Table 4.} \emph{Descriptive Statistics for Study Variables}

% Table shown without inline diff markup for longtable stability.
\begin{longtable}[]{@{}llllll@{}}
\toprule
Variable & \emph{n} & \emph{M} & \emph{SD} & Min & Max \\
\midrule
\endhead
\bottomrule
\endlastfoot
Present-oriented perspective (reversed) & 287 & 2.97 & 0.64 & 1.33 & 4.00 \\
Role of the historical agent & 287 & 2.80 & 0.68 & 1.00 & 4.00 \\
Historical contextualization & 287 & 2.71 & 0.74 & 1.00 & 4.00 \\
Six-item contextualization composite & 287 & 2.84 & 0.55 & 1.33 & 4.00 \\
Nine-item HPT total score & 287 & 2.83 & 0.49 & 1.67 & 3.89 \\
Historical knowledge (0--6) & 293 & 3.04 & 1.62 & 0.00 & 6.00 \\
Dictatorship acceptance & 285 & 2.54 & 0.88 & 1.00 & 5.00 \\
Nazi-crime relativization & 286 & 2.43 & 0.89 & 1.00 & 5.00 \\
Right-authoritarian attitudes & 284 & 2.48 & 0.75 & 1.00 & 5.00 \\
Authoritarianism & 283 & 2.86 & 0.62 & 1.00 & 4.60 \\
Social desirability & 283 & 3.02 & 0.63 & 1.00 & 4.60 \\
\end{longtable}

\emph{Note.} HPT items scored on a 1--4 scale; POP items reverse-coded. FR-LF and KSA-3 scored on a 1--5 scale. KN scored as sum of correct responses (0--6).

\hypertarget{intraclass-correlations}{%
\subsubsection{4.5 Intraclass Correlations}\label{intraclass-correlations}}

Intraclass correlations indicated meaningful school-level clustering for most HPT outcomes (\DIFdelbegin \DIFdel{HPT\_CTX6}\DIFdelend \DIFaddbegin \DIFadd{six-item contextualization composite}\DIFaddend : ICC\_school = .041; \DIFdelbegin \DIFdel{HPT\_TOT9}\DIFdelend \DIFaddbegin \DIFadd{nine-item total score}\DIFaddend : ICC\_school = .085; \DIFdelbegin \DIFdel{ROA}\DIFdelend \DIFaddbegin \DIFadd{role of the historical agent}\DIFaddend : ICC\_school = .102), consistent with between-school differences in instructional emphasis and student composition. Classroom-level variance was near zero for the primary composite but non-negligible for \DIFdelbegin \DIFdel{CONT }\DIFdelend \DIFaddbegin \DIFadd{contextualization }\DIFaddend (ICC\_class = .033) and \DIFdelbegin \DIFdel{ROA }\DIFdelend \DIFaddbegin \DIFadd{role of the historical agent }\DIFaddend (ICC\_class = .104). These ICCs justify multilevel modelling for substantive analyses.
The school-level variance is consistent with between-school
differences in history instruction, curricular emphasis, and student
composition\DIFdelbegin \DIFdel{. The corresponding design effects for the primary composite were modest (HPT\_CTX6: DEFF = 1.56; HPT\_TOT9: DEFF = 2.16; see Section 3.4 for computation and implications)}\DIFdelend \DIFaddbegin \DIFadd{, although the present design cannot identify which of
these mechanisms produced the variation}\DIFaddend .

\hypertarget{dif-analysis}{%
\subsubsection{4.6 DIF Analysis}\label{dif-analysis}}

Table 5 presents \DIFaddbegin \DIFadd{omnibus }\DIFaddend DIF results for all nine HPT items. Using a constrained multi-group graded response model with likelihood-ratio tests and Bonferroni-adjusted significance (\(\alpha\) = .01), no items were flagged\DIFdelbegin \DIFdel{for either uniform or non-uniform DIF}\DIFdelend . All adjusted \emph{p}-values exceeded .01\DIFdelbegin \DIFdel{, indicating statistically equivalent item parameters across ideology groups. This means that at equal levels of underlying HPT competence, students with low and high ideological attitudes responded to each item in the same way---no individual item functioned differentially as a function of political attitudes}\DIFdelend . \DIFaddbegin \DIFadd{No DIF was detected between the low- and high-ideology groups with this sample and procedure; the magnitude that the design could reliably detect is unknown.
}\DIFaddend

\textbf{Table 5.} \emph{DIF Results per HPT Item (Likelihood-Ratio Tests from Multi-Group Graded Response Model)}

% Table shown without inline diff markup for longtable stability.
\begin{longtable}[]{@{}llllll@{}}
\toprule
Item & \(\chi^2\) & \emph{df} & \emph{p} (raw) & \emph{p} (Bonferroni) & Flagged \\
\midrule
\endhead
\bottomrule
\endlastfoot
POP1 & 10.46 & 4 & .033 & .300 & No \\
POP2 & 9.01 & 4 & .061 & .547 & No \\
POP3 & 9.72 & 4 & .045 & .408 & No \\
ROA1 & 9.59 & 4 & .048 & .432 & No \\
ROA2 & 3.13 & 4 & .536 & 1.000 & No \\
ROA3 & 8.25 & 4 & .083 & .744 & No \\
CONT1 & 1.51 & 4 & .825 & 1.000 & No \\
CONT2 & 3.60 & 4 & .462 & 1.000 & No \\
CONT3 & 9.63 & 4 & .047 & .424 & No \\
\end{longtable}

\emph{Note.} Likelihood-ratio chi-square statistics from constrained
multi-group graded response models. Each test compares a fully
constrained baseline (all items equal across groups) to a model
freeing discrimination and \DIFdelbegin \DIFdel{threshold }\DIFdelend \DIFaddbegin \DIFadd{category-intercept }\DIFaddend parameters for the tested item.
Low (\emph{n} = 96) and High (\emph{n} = \DIFdelbegin \DIFdel{96}\DIFdelend \DIFaddbegin \DIFadd{95}\DIFaddend ) ideology groups formed
by tertile split on composite ideology score (middle tertile
excluded). Bonferroni-adjusted \(\alpha\) = .01 (9 tests). 0 of 9
items flagged. POP1 showed the largest DIF statistic (\(\chi^2\) =
\DIFdelbegin \DIFdel{11.18}\DIFdelend \DIFaddbegin \DIFadd{10.46}\DIFaddend , \emph{p}\_raw = \DIFdelbegin \DIFdel{.025}\DIFdelend \DIFaddbegin \DIFadd{.033}\DIFaddend ), well above the Bonferroni threshold
after correction (\emph{p}\_adj = \DIFdelbegin \DIFdel{.222}\DIFdelend \DIFaddbegin \DIFadd{.300}\DIFaddend ). GRM item discrimination and
\DIFdelbegin \DIFdel{threshold }\DIFdelend \DIFaddbegin \DIFadd{category-intercept }\DIFaddend parameters are reported in Table S1 in the OSF
supplementary materials.

\DIFdelbegin \DIFdel{Although no individual item reached significance after Bonferroni
correction, five of nine items (POP1}\DIFdelend \DIFaddbegin \DIFadd{Because every other item served as an anchor}\DIFaddend , \DIFdelbegin \DIFdel{POP3, ROA1, ROA3, CONT3) had
raw }\emph{\DIFdel{p}}%DIFAUXCMD
\DIFdel{-values below .10. This clustering of near-significant
results is consistent with a pattern in which multiple items share
weak DIF in the same direction---precisely the scenario under which
the all-others-as-anchor strategy without iterative purification is
most vulnerable to anchorcontamination }%DIFDELCMD < \citep{Woods2009}%%%
\DIFdel{. Because the
anchor set may itself be contaminated by small uniform DIF, the
likelihood-ratio tests may underestimate the true magnitude of
item-level differences. This pattern does not invalidate the overall
null DIF finding, but it suggests that the absence of flagged items
should }\DIFdelend \DIFaddbegin \DIFadd{undetected DIF among
those anchors could affect each test. The result should therefore }\DIFaddend be
interpreted as ``no \DIFdelbegin \DIFdel{moderate-to-large DIF detected '' rather
than }\DIFdelend \DIFaddbegin \DIFadd{DIF detected with this sample and procedure,'' not
}\DIFaddend ``no DIF present.''

\hypertarget{measurement-invariance}{%
\subsubsection{4.7 Measurement Invariance}\label{measurement-invariance}}

Table 6 presents the stepwise MG-CFA results. The configural model \DIFdelbegin \DIFdel{demonstrated adequate fit (}\DIFdelend \DIFaddbegin \DIFadd{produced }\DIFaddend CFI = \DIFdelbegin \DIFdel{.978}\DIFdelend \DIFaddbegin \DIFadd{.977}\DIFaddend , RMSEA = \DIFdelbegin \DIFdel{.043, }\DIFdelend \DIFaddbegin \DIFadd{.044, and }\DIFaddend SRMR = .074\DIFdelbegin \DIFdel{), confirming that the three-factor structure holds separately in both ideology groups. Constraining factor loadings to equality across groups (metric invariance) produced no meaningful deterioration (}\DIFdelend \DIFaddbegin \DIFadd{. Adding equality constraints did not worsen global fit: metric }\DIFaddend \(\Delta\)CFI = \DIFdelbegin \DIFdel{.000, }\DIFdelend \DIFaddbegin \DIFadd{$-$.001 and }\DIFaddend \(\Delta\)RMSEA = $-$\DIFdelbegin \DIFdel{.002). Constraining both loadings and thresholds (scalar invariance) similarly maintained fit (}\DIFdelend \DIFaddbegin \DIFadd{.001; scalar }\DIFaddend \(\Delta\)CFI = +\DIFdelbegin \DIFdel{.002, }\DIFdelend \DIFaddbegin \DIFadd{.005 and }\DIFaddend \(\Delta\)RMSEA = $-$\DIFdelbegin \DIFdel{.006). Both transitions satisfy the established criteria of $|\Delta\text{CFI}|$ \(\leq\) .01 and $|\Delta\text{RMSEA}|$ \(\leq\) .015 }\DIFdelend \DIFaddbegin \DIFadd{.008. These changes fall within common fit-index criteria }\DIFaddend \citep{Chen2007,Cheung2002}\DIFdelbegin \DIFdel{, providing evidence for full scalar invariance. }\DIFdelend \DIFaddbegin \DIFadd{. All three models converged and passed lavaan's post-estimation admissibility check; the diagonals of their model-implied residual-variance matrices were positive.
}\DIFaddend

\textbf{Table 6.} \emph{Multi-Group CFA Fit Indices Across Invariance Levels (WLSMV)}

% Table shown without inline diff markup for longtable stability.
\begin{longtable}[]{@{}llllll@{}}
\toprule
Model & CFI & RMSEA & SRMR & \(\Delta\)CFI & \(\Delta\)RMSEA \\
\midrule
\endhead
\bottomrule
\endlastfoot
Configural & .977 & .044 & .074 & --- & --- \\
Metric & .976 & .043 & .081 & $-$.001 & $-$.001 \\
Scalar & .981 & .035 & .076 & +.005 & $-$.008 \\
\end{longtable}

\emph{Note.} \(\Delta\)CFI and \(\Delta\)RMSEA computed relative to the preceding model. Invariance criteria: $|\Delta\text{CFI}|$ \(\leq\) .01, $|\Delta\text{RMSEA}|$ \(\leq\) .015 \citep{Chen2007}. Low and High ideology groups formed by tertile split on composite ideology\DIFdelbegin \DIFdel{(}\emph{\DIFdel{n}} %DIFAUXCMD
\DIFdel{\(\approx\) 96 per group)}\DIFdelend \DIFaddbegin \DIFadd{; complete item-response data were available for 92 Low and 90 High participants}\DIFaddend . Items treated as ordered categorical.

The SRMR increased slightly from configural (.074) to metric (.081),
marginally exceeding the conventional \(\leq\) .08 threshold. This
threshold was developed for maximum-likelihood estimation with
continuous indicators and may not apply directly to WLSMV-estimated
models with polychoric residuals; the marginal exceedance should
therefore be interpreted cautiously. The SRMR returned to .076 at the
scalar step. \DIFdelbegin \DIFdel{This pattern can signal partial loading non-invariance.
The largest }\DIFdelend \DIFaddbegin \DIFadd{The largest ordinary }\DIFaddend modification index at the metric step
was \DIFdelbegin \DIFdel{11.28 (a
suggested cross-loading of }\DIFdelend \DIFaddbegin \DIFadd{11.22 and suggested a }\DIFaddend POP2 \DIFaddbegin \DIFadd{cross-loading }\DIFaddend on the ROA factor in the
Low \DIFdelbegin \DIFdel{ideology
group), marginally above the conventional threshold of
10. This
suggests a minor loading differencefor POP2 that is consistent with
its weak loading (.30) in the overall CFA (Table 2) rather than
evidence of systematic group-specific misspecification}\DIFdelend \DIFaddbegin \DIFadd{group. Because ordinary modification indices ignore equality
constraints, this is a local cross-loading diagnostic, not evidence of
a between-group loading difference}\DIFaddend .

\DIFdelbegin \DIFdel{Negative estimated residual variances (Heywood cases) occurred for ROA
items in the Low ideology group in the configural and metric models.
Such estimates indicate a non-positive-definite covariance matrix,
which can invalidate fit statistics and render parameter estimates
uninterpretable. While Heywood cases can arise with small per-group
samples and WLSMV estimation of ordinal items---particularly with
three-item factors---they signal potential modelinadmissibility }\DIFdelend \DIFaddbegin \DIFadd{The admissible solutions and changes in CFI and RMSEA are compatible
with scalar invariance under the fitted model. This remains qualified
evidence }\DIFaddend rather than a \DIFdelbegin \DIFdel{benign estimation artefact. The per-group sample sizes
(}\emph{\DIFdel{n}} %DIFAUXCMD
\DIFdel{\(\approx\) 82--96) are at the lower boundary of what
MG-CFA with WLSMV requires }%DIFDELCMD < \citep{Putnick2016}%%%
\DIFdel{. The
configural and metric invariance conclusions should therefore be
qualified: they represent the best available evidence given sample size
constraints, but the Heywood cases introduce uncertainty about the
precision of the configural baseline.
The invariance results should be
treated as convergent evidence alongside the DIF analysis rather than
as definitive confirmation of scalar equivalence.
}\DIFdelend \DIFaddbegin \DIFadd{firm equivalence conclusion: the analysis was
added after registration, used only 90--92 complete cases per extreme
group, estimated three-indicator factors, and did not model school or
classroom clustering.
}\DIFaddend

\hypertarget{supplementary-invariance-ns-only-grouping}{%
\subsubsection{4.8 Supplementary Invariance: NS-Only Grouping}\label{supplementary-invariance-ns-only-grouping}}

To verify that invariance findings were not an artefact of the
composite ideology grouping variable, we repeated the MG-CFA using
groups defined by the NS facet alone (National Socialist
sympathy/\DIFdelbegin \DIFdel{relativisation}\DIFdelend \DIFaddbegin \DIFadd{relativization}\DIFaddend ), which represents the most direct theoretical
\DIFdelbegin \DIFdel{operationalisation }\DIFdelend \DIFaddbegin \DIFadd{operationalization }\DIFaddend of content congruence with the Weimar scenario. With
the NS tertile split, group sizes were \emph{n} = 126 (Low) and
\emph{n} = \DIFdelbegin \DIFdel{125 }\DIFdelend \DIFaddbegin \DIFadd{123 }\DIFaddend (High)\DIFdelbegin \DIFdel{. The supplementary analysis largely supported
scalar invariance }\DIFdelend \DIFaddbegin \DIFadd{, with 37 middle-tertile cases and 7 missing NS scores
excluded. Fit did not
deteriorate when equality constraints were added }\DIFaddend (configural CFI = \DIFdelbegin \DIFdel{.960}\DIFdelend \DIFaddbegin \DIFadd{.959}\DIFaddend ,
RMSEA = \DIFdelbegin \DIFdel{.050}\DIFdelend \DIFaddbegin \DIFadd{.051}\DIFaddend , SRMR = \DIFdelbegin \DIFdel{.070}\DIFdelend \DIFaddbegin \DIFadd{.071}\DIFaddend ;
metric \(\Delta\)CFI = +.002, \(\Delta\)SRMR = +.006; scalar
\(\Delta\)CFI = +.014, \(\Delta\)SRMR = $-$.003). \DIFdelbegin \DIFdel{The scalar
\(\Delta\)CFI
of .014 marginally exceeds the conventional \(\leq\)
.01 threshold }%DIFDELCMD < \citep{Chen2007}%%%
\DIFdel{, indicating that full scalar invariance is
not unambiguously confirmed for the NS-only grouping. The
}\DIFdelend \DIFaddbegin \DIFadd{Because CFI
increased and RMSEA and SRMR decreased at the scalar step, the positive
}\DIFaddend \(\Delta\)\DIFdelbegin \DIFdel{RMSEA remained within acceptable bounds, and the direction
of CFI change was positive
(improved fit ), which is atypical for
genuine invariance violations. Nevertheless, given that the NS facet
is the most theoretically proximal grouping variable and the larger
per-group sample sizes (}\emph{\DIFdel{n}} %DIFAUXCMD
\DIFdel{= 126/125) make the criterion breach
more interpretable, partial scalar invariance---with possible
threshold non-equivalence for individual items---cannot be ruled out. }\DIFdelend \DIFaddbegin \DIFadd{CFI does not indicate worse fit and should not be compared
with a deterioration threshold by absolute magnitude. This
post-registration analysis provides descriptive robustness evidence,
not a second confirmatory test. }\DIFaddend Full fit statistics for this
supplementary analysis are available in the OSF repository.

\hypertarget{convergent-validity}{%
\subsubsection{4.9 Convergent Validity}\label{convergent-validity}}

Scenario-relevant historical knowledge correlated positively with all
HPT outcomes: \emph{r} = .34 \DIFdelbegin \DIFdel{(HPT\_CTX6)}\DIFdelend \DIFaddbegin \DIFadd{for the six-item contextualization
composite}\DIFaddend , \emph{r} = .22 \DIFdelbegin \DIFdel{(CONT), }\DIFdelend \DIFaddbegin \DIFadd{for contextualization, and }\DIFaddend \emph{r} = .32
\DIFdelbegin \DIFdel{(POP\_rev)}\DIFdelend \DIFaddbegin \DIFadd{for reversed present-oriented perspective}\DIFaddend . These moderate correlations indicate that
students with stronger knowledge of the relevant historical period
scored higher on HPT\DIFdelbegin \DIFdel{, while also demonstrating that HPT is not reducible to knowledge alone---the shared variance (approximately
5--12\%) leaves substantial HPT variance attributable to
reasoning processes beyond factual recall}\DIFdelend \DIFaddbegin \DIFadd{. The scores were not empirically interchangeable,
but their unshared variance cannot be attributed specifically to
reasoning because it also includes measurement error and other
unmeasured influences}\DIFaddend . A limitation of this evidence is that
the KN items were written to cover content directly relevant to the
HPT scenario, so the correlation may partly reflect shared
specific-content activation rather than the broader
``knowledge-dependent competence'' relationship posited theoretically.
A general historical knowledge measure would provide a stronger
convergent validity test.

History grade also correlated with \DIFdelbegin \DIFdel{HPT\_CTX6 }\DIFdelend \DIFaddbegin \DIFadd{the six-item contextualization composite }\DIFaddend (\emph{r} = $-$.153, \emph{p} = .010; negative sign reflects the Czech scale where 1 = excellent). By contrast, \DIFdelbegin \DIFdel{ideology measures showed negligible correlations with HPT (FR-LF--HPT\_CTX6: }\DIFdelend \DIFaddbegin \DIFadd{right-authoritarian attitudes showed a negligible correlation with the composite (}\DIFaddend \emph{r} = $-$.05), and social desirability was unrelated \DIFaddbegin \DIFadd{to HPT outcomes }\DIFaddend (\emph{r} = $-$.02 to .01). \DIFdelbegin \DIFdel{This pattern---significant knowledge correlations, null ideology and social desirability correlations---is consistent with an instrument that measures disciplinary reasoning rather than attitudes or impression management. }\DIFdelend \DIFaddbegin \DIFadd{One exploratory exception was a negative association between authoritarianism and reversed present-oriented perspective (}\emph{\DIFadd{r}} \DIFadd{= $-$.13, 95\% CI }[\DIFadd{$-$.25, $-$.02}]\DIFadd{). The corresponding mutually adjusted standardized KSA-3 coefficient was \(\beta\) = $-$.166, }\emph{\DIFadd{SE}} \DIFadd{= .066, unadjusted }\emph{\DIFadd{p}} \DIFadd{= .013. This association was not preregistered or subjected to the preregistered multiplicity procedure and should be treated as a hypothesis-generating pattern consistent with the reduced-integrative-complexity pathway rather than a confirmed effect. Overall, the positive knowledge associations and absence of detected associations for the primary ideology--HPT relationship are consistent with, but do not establish, the intended score interpretation.
}\DIFaddend

\textbf{Table 7.} \emph{Zero-Order Pearson Correlations Among Key Variables}

\DIFaddbegin \begin{scriptsize}
\DIFaddend % Table shown without inline diff markup for longtable stability.
\begin{longtable}[]{@{}p{2.7cm}rrrrrrr@{}}
\toprule
 & HPT & Context & \shortstack{Presentism\\(reversed)} & Knowledge & \shortstack{Right-authoritarian\\attitudes} & Authoritarianism & \shortstack{Social\\desirability} \\
\midrule
\endhead
\bottomrule
\endlastfoot
HPT composite & --- & [.78, .86] & [.71, .81] & [.23, .43] & [$-$.16, .07] & [$-$.15, .09] & [$-$.14, .10] \\
Contextualization & .82*** & --- & [.14, .36] & [.11, .33] & [$-$.10, .13] & [$-$.05, .19] & [$-$.11, .13] \\
Present-oriented perspective (rev.) & .76*** & .26*** & --- & [.20, .41] & [$-$.21, .02] & [$-$.25, $-$.02] & [$-$.16, .07] \\
Knowledge & .34*** & .23*** & .31*** & --- & [$-$.22, .01] & [$-$.07, .16] & [$-$.09, .15] \\
Right-authoritarian attitudes & $-$.05 & .02 & $-$.10 & $-$.10 & --- & [.43, .60] & [$-$.30, $-$.08] \\
Authoritarianism & $-$.03 & .07 & $-$.13* & .05 & .52*** & --- & [$-$.26, $-$.03] \\
Social desirability & $-$.02 & .01 & $-$.05 & .03 & $-$.19** & $-$.15* & --- \\
\end{longtable}
\DIFaddbegin \end{scriptsize}
\DIFaddend

\emph{Note.} \DIFaddbegin \DIFadd{Correlations are below the diagonal and unadjusted descriptive 95\% confidence intervals are above it. }\DIFaddend Pairwise complete observations (\emph{n} = \DIFdelbegin \DIFdel{276--287). HPT POP }\DIFdelend \DIFaddbegin \DIFadd{282--287; exact values in Table S3b). Present-oriented perspective }\DIFaddend is reverse-scored. \DIFdelbegin \DIFdel{KN }\DIFdelend \DIFaddbegin \DIFadd{Knowledge }\DIFaddend = historical knowledge (0--6 sum); FR-LF = right-authoritarian ideology (1--5); KSA-3 = authoritarianism (1--5); SDR-5 = social desirability (1--5). \DIFaddbegin \DIFadd{Intervals and }\emph{\DIFadd{p}}\DIFadd{-values do not adjust for school or classroom clustering. }\DIFaddend *** \emph{p} $<$ .001. ** \emph{p} $<$ .01. * \emph{p} $<$ .05.

\DIFaddbegin \DIFadd{Figure~\ref{fig:relationships} visualizes the near-zero ideology
association and positive knowledge association with the six-item
contextualization composite.
}

\begin{figure}[htbp]
\centering
\includegraphics[width=\textwidth]{revision_outputs/revision_relationships.pdf}
\caption{\DIFaddFL{Observed relationships of the six-item contextualization composite with the combined ideology score---the mean of standardized right-authoritarian attitude and authoritarianism scores---(Panel A) and historical knowledge (Panel B). Points are jittered to display overlap; lines and shaded regions show unadjusted linear fits and 95\% confidence intervals.}}
\label{fig:relationships}
\end{figure}

\DIFaddend \hypertarget{substantive-context-ideology-hpt-relationships}{%
\subsubsection{4.10 Substantive Context: Ideology-HPT Relationships}\label{substantive-context-ideology-hpt-relationships}}

\DIFdelbegin \DIFdel{Ideology was unrelated to HPT scores across all model
specifications: FR-LF }\DIFdelend \DIFaddbegin \DIFadd{The mutually adjusted standardized right-authoritarian attitude coefficient
was near zero for the primary six-item
contextualization composite: }\DIFaddend \(\beta\) = \DIFdelbegin \DIFdel{0.012}\DIFdelend \DIFaddbegin \DIFadd{.016}\DIFaddend , \emph{SE} = \DIFdelbegin \DIFdel{0.067}\DIFdelend \DIFaddbegin \DIFadd{.068}\DIFaddend , 95\%
CI [$-$\DIFdelbegin \DIFdel{0.120, 0.143}\DIFdelend \DIFaddbegin \DIFadd{.117, .149}\DIFaddend ], \emph{p} = \DIFdelbegin \DIFdel{.863 for the primary outcome
(HPT\_CTX6).
}\DIFdelend \DIFaddbegin \DIFadd{.817 (analysis }\emph{\DIFadd{n}} \DIFadd{= 282).
The fitted model issued a singularity warning because the classroom
random-intercept variance was estimated as zero; the school intercept and
fixed effects remained estimable, and school fixed-effects sensitivity models
gave the same substantive conclusion.
}\DIFaddend Historical knowledge was the only consistent
predictor (\(\beta\) = 0.21--0.34, \emph{p} $<$ .001).
As an exploratory sensitivity analysis, equivalence tests used the Two
One-Sided Tests (TOST) procedure, which reverses the logic of
null-hypothesis testing by asking whether an effect is small enough to
be declared practically negligible. The Smallest Effect Size Of
Interest (SESOI) was set post hoc at \(\pm\) 0.20 \DIFdelbegin \DIFdel{standardised }\DIFdelend \DIFaddbegin \DIFadd{standardized }\DIFaddend units,
corresponding to \citeauthor{Cohen1988}'s (\citeyear{Cohen1988}) threshold for a ``small'' effect.
This SESOI is wider than the preregistered power target of \(\beta\)
\(\geq\) .15, meaning that effects between .15 and .20 cannot be
confidently excluded; the TOST results should therefore be treated as
exploratory rather than confirmatory. With the \(\pm\) 0.20 boundary,
the TOST results indicated that \DIFdelbegin \DIFdel{ideology }\DIFdelend \DIFaddbegin \DIFadd{right-authoritarian attitude }\DIFaddend effects at
or above \(\beta\) = .20 are statistically ruled out \DIFdelbegin \DIFdel{(all }\DIFdelend \DIFaddbegin \DIFadd{for the primary
relationship (}\DIFaddend TOST \emph{p} \DIFdelbegin \DIFdel{s $<$ .003}\DIFdelend \DIFaddbegin \DIFadd{= .0034}\DIFaddend ; 90\% CI \DIFdelbegin \DIFdel{for the primary outcome }\DIFdelend [$-$\DIFdelbegin \DIFdel{0.10, 0.12}\DIFdelend \DIFaddbegin \DIFadd{.096, .127}\DIFaddend ], entirely
within the \(\pm\) 0.20 SESOI). The \DIFdelbegin \DIFdel{ideology coefficient }\DIFdelend \DIFaddbegin \DIFadd{right-authoritarian attitude
coefficient for the primary composite }\DIFaddend remained non-significant
across all sensitivity checks, including alternative HPT scoring
composites, different ideology operationalizations, exclusion of
high social desirability responders, random-slope models, and
school fixed-effects specifications\DIFdelbegin \DIFdel{(all }\emph{\DIFdel{p}}%DIFAUXCMD
\DIFdel{s $>$
.43)}\DIFdelend . These substantive findings complement the measurement
evidence: \DIFdelbegin \DIFdel{the instrument not only measures equivalently across
ideology groups (DIF , invariance) but also yields scores that are
uncontaminated by political attitudes at the composite level}\DIFdelend \DIFaddbegin \DIFadd{no DIF or composite-level ideology
association was detected within the observed ideological range}\DIFaddend .

\hypertarget{discussion}{%
\subsection{5. Discussion}\label{discussion}}

\hypertarget{summary-of-findings}{%
\subsubsection{5.1 Summary of Findings}\label{summary-of-findings}}

The \DIFdelbegin \DIFdel{HPT instrument passed every fairness test applied. Factor
structure replicated, DIF was absent, scalar invariance held, and
ideology }\DIFdelend \DIFaddbegin \DIFadd{results provide qualified fairness evidence. The expected factor
structure was supported, no DIF was detected with this procedure, and
right-authoritarian attitudes }\DIFaddend did not predict \DIFdelbegin \DIFdel{HPT scores. This is the first psychometric evaluation of a Czech adaptation of the Hartmann and Hasselhorn HPT instrument and the first investigation of ideological fairness for any HPT measure}\DIFdelend \DIFaddbegin \DIFadd{the primary HPT composite;
fit-index changes were compatible with scalar invariance, subject to small
extreme groups and unmodelled clustering}\DIFaddend . The Czech adaptation \DIFdelbegin \DIFdel{demonstrated the theorized three-factor structure, acceptable reliabilityfor three-item subscales, meaningful convergent validity with }\DIFdelend \DIFaddbegin \DIFadd{showed modest subscale reliability, a positive relationship with scenario-relevant }\DIFaddend historical knowledge, and sensitivity to school-level variation. \DIFdelbegin \DIFdel{Historical knowledge was the only construct that predicted HPT scores; political ideology , authoritarianism, and social desirability were all unrelated}\DIFdelend \DIFaddbegin \DIFadd{An unadjusted exploratory authoritarianism--reversed-presentism association was not subjected to the preregistered multiplicity procedure; other ideology and social-desirability associations were not detected, although weak reliability constrains the latter null associations}\DIFaddend .

\DIFdelbegin %DIFDELCMD < \hypertarget{factor-structure-confirms-the-three-level-hpt-model}{%
%DIFDELCMD < \subsubsection{5.2 Factor Structure Confirms the Three-Level HPT Model}\label{factor-structure-confirms-the-three-level-hpt-model}}
%DIFDELCMD < %%%
\DIFdelend \DIFaddbegin \hypertarget{factor-structure-evidence}{%
\subsubsection{5.2 Factor-Structure Evidence}\label{factor-structure-evidence}}
\DIFaddend

The \DIFaddbegin \DIFadd{correlated }\DIFaddend three-factor model (POP, ROA, CONT) fit \DIFdelbegin \DIFdel{substantially better than alternatives, confirming that the Czech adaptation preserves the dimensional structure of }\DIFdelend \DIFaddbegin \DIFadd{the Czech data better than the tested alternatives. This does not demonstrate cross-cultural equivalence, and the specification differs from the two-dimensional PCA pattern in }\DIFaddend the original German \DIFdelbegin \DIFdel{instrument. This replication in a linguistically and educationally distinct context confirms that the item grouping holds in a Czech sample }%DIFDELCMD < \citep{Hambleton2005,ITC2017}%%%
\DIFdel{, though with }\DIFdelend \DIFaddbegin \DIFadd{study. It shows only that the Czech data are compatible with separating the three item sets, under }\DIFaddend a near-saturated model \DIFdelbegin \DIFdel{the structural confirmation is relatively undemanding (see Section 4.1)}\DIFdelend \DIFaddbegin \DIFadd{with limited opportunity to reveal misfit}\DIFaddend . CONT items loaded most consistently (.61--.72), while POP2 and ROA2 showed weaker loadings (.30); these were retained for comparability but are candidates for revision in future adaptations.

\hypertarget{addressing-modest-reliability}{%
\subsubsection{5.3 Addressing Modest Reliability}\label{addressing-modest-reliability}}

The subscale reliabilities (\(\alpha\) = \DIFdelbegin \DIFdel{.47}\DIFdelend \DIFaddbegin \DIFadd{.46}\DIFaddend --.64; \(\omega\) = \DIFdelbegin \DIFdel{.55}\DIFdelend \DIFaddbegin \DIFadd{.53}\DIFaddend --\DIFdelbegin \DIFdel{.69}\DIFdelend \DIFaddbegin \DIFadd{.70}\DIFaddend ) fall below conventional benchmarks \DIFdelbegin \DIFdel{but are structurally constrained by the three-item design---a feature of the original instrument, not a deficiency of the Czech adaptation. With }\emph{\DIFdel{k}} %DIFAUXCMD
\DIFdelend \DIFaddbegin \DIFadd{and limit individual-level interpretation. Three-item scales make high coefficients difficult to attain, but brevity does not make the observed precision adequate. }\citet{Loevinger1957Objective} \DIFadd{cautioned that highly homogeneous item sets can reduce a construct to repeated versions of a narrow question and argued for broad content representation; that principle supports balancing coverage and homogeneity, not treating low reliability as harmless. The six- and nine-item HPT scores were more precise (\(\omega\) }\DIFaddend = \DIFdelbegin \DIFdel{3 }\DIFdelend \DIFaddbegin \DIFadd{.66 }\DIFaddend and \DIFdelbegin \DIFdel{inter-item correlations of .22--.37}\DIFdelend \DIFaddbegin \DIFadd{.70)}\DIFaddend , \DIFdelbegin \DIFdel{Spearman-Brown projections yield \(\alpha\) = .64--.79 for six-item subscales. Despite modest reliability, the instrument detected theoretically expected knowledge--HPT relationships (\(r\) = .22--.34) and discriminated between HPT levels in the expected ordering. The disattenuated ideology--HPT correlation (}\emph{\DIFdel{r}} %DIFAUXCMD
\DIFdel{= -.065)confirms that even with perfect reliability, no meaningful ideology association would emerge. We recommend using composite scores (HPT\_CTX6 or HPT\_TOT9) }\DIFdelend \DIFaddbegin \DIFadd{so we recommend composite scores }\DIFaddend rather than individual subscales for \DIFdelbegin \DIFdel{applied purposes. }\DIFdelend \DIFaddbegin \DIFadd{group-level research. Associations involving the least reliable measures, especially SDR-5, should be read cautiously because attenuation may obscure weak relationships.
}\DIFaddend

\hypertarget{measurement-fairness-evidence}{%
\subsubsection{5.4 Measurement Fairness Evidence}\label{measurement-fairness-evidence}}

\DIFdelbegin \DIFdel{Scalar invariance was confirmed for the composite ideology grouping
and largely supported for the NS-only analysis (with the
caveat that
the NS scalar \(\Delta\)CFI marginally exceeded the conventional
threshold; see Section 4.8). These results indicate that the HPT
instrument imposes largely equivalent measurement demands on students
regardless of their political attitudes}\DIFdelend \DIFaddbegin \DIFadd{Fit-index changes were compatible with scalar invariance for both
ideology groupings, and the composite-group models were admissible in the
current reproducible analysis. The small extreme groups, unmodelled
clustering, and post-registration status nevertheless prevent a firm
equivalence conclusion}\DIFaddend . The DIF analysis
\DIFdelbegin \DIFdel{complements
this at the item level: no item exhibited differential functioning
after Bonferroni correction, which is particularly important for the CONT items most theoretically susceptible to ideological
contamination, though the clustering of near-significant raw
}\emph{\DIFdel{p}}%DIFAUXCMD
\DIFdel{-values warrants caution (see Section 4.6).
}\DIFdelend \DIFaddbegin \DIFadd{adds item-level evidence, but no design-specific power analysis was
conducted and the unidimensional GRM does not match the CFA structure.
Together, these results provide qualified rather than definitive
fairness evidence.
}\DIFaddend

The correlation pattern---positive scenario-relevant knowledge--HPT
associations \DIFdelbegin %DIFDELCMD < \citep{Chapman2021Introduction,Hartmann2008,Wineburg2001Historical}%%%
\DIFdel{, null ideology correlations, }\DIFdelend \DIFaddbegin \citep{Chapman2021Introduction,Hartmann_Hasselhorn2008Historical,Wineburg2001Historical}\DIFadd{, a near-zero primary ideology correlation, and }\DIFaddend null social desirability
correlations---is consistent with the intended score interpretation
\citep{AERA2014Standards,Messick1995Validity}. However, the convergent evidence is bounded by the
scenario-proximal nature of the knowledge measure (see Section 4.9),
which limits the strength of the knowledge-dependent competence claim.

The near-zero ideology--HPT correlation (\emph{r} = $-$.05) should
not be interpreted as conclusive evidence that ideology plays no role
in HPT performance. As discussed in Section 2.3, two opposing
pathways are theoretically plausible: a congruence pathway (whereby
attitudinal alignment inflates scores) and a reduced integrative
complexity pathway (whereby authoritarian cognition depresses
\DIFdelbegin \DIFdel{contextualisation}\DIFdelend \DIFaddbegin \DIFadd{contextualization}\DIFaddend ). If both pathways operate simultaneously, their
effects could cancel, producing a near-zero net correlation that masks
two moderate-sized opposing effects. The present data cannot
distinguish a genuine null from cancelling pathways.

The present fairness evidence applies to mainstream Czech adolescent
samples with restricted ideological range. Users in contexts with
greater political \DIFdelbegin \DIFdel{polarisation }\DIFdelend \DIFaddbegin \DIFadd{polarization }\DIFaddend should conduct their own DIF and
invariance analyses before assuming equivalent measurement.

\hypertarget{implications-for-assessment-practice}{%
\subsubsection{5.5 Implications for Assessment Practice}\label{implications-for-assessment-practice}}

The core practical message is that politically charged assessment content \DIFdelbegin \DIFdel{does not automatically introduce }\DIFdelend \DIFaddbegin \DIFadd{need
not produce detectable }\DIFaddend construct-irrelevant \DIFdelbegin \DIFdel{variance---but this conclusion }\DIFdelend \DIFaddbegin \DIFadd{variance in every application. In
this study, no primary composite association or DIF flag was detected at the
available sensitivity. This pattern }\DIFaddend is bounded by the \DIFdelbegin \DIFdel{present sample's }\DIFdelend restricted ideological
range and should not be generalized beyond \DIFaddbegin \DIFadd{the single }\DIFaddend structured,
closed-ended \DIFdelbegin \DIFdel{formats }\DIFdelend \DIFaddbegin \DIFadd{scenario }\DIFaddend without further evidence.

The evaluation protocol used here \DIFdelbegin \DIFdel{provides an evidence-based model }\DIFdelend \DIFaddbegin \DIFadd{illustrates a possible workflow }\DIFaddend for
evaluating fairness when assessment content intersects with respondents'
attitudes. \DIFdelbegin \DIFdel{This protocol is not specific to HPTor history education; it applies to any }\DIFdelend \DIFaddbegin \DIFadd{Although developed for HPT, the workflow may inform studies of
other }\DIFaddend scenario-based assessment\DIFaddbegin \DIFadd{s}\DIFaddend  where content carries attitudinal \DIFdelbegin \DIFdel{valence---moral reasoning instruments using politically divisive dilemmas, science literacy assessments touching evolution or climate change, and civic education assessments involving controversial policies}\DIFdelend \DIFaddbegin \DIFadd{valence}\DIFaddend .
The general approach involves five steps:

\begin{enumerate}
\def\labelenumi{\arabic{enumi}.}
\tightlist
\item
  Identify the attitudinal dimension theoretically most likely to produce CIV.
\item
  Form extreme groups on that dimension.
\item
  Test DIF at the item level.
\item
  Test measurement invariance at the model level.
\item
  Examine \DIFdelbegin \DIFdel{convergent and discriminant correlations to verify that construct scores relate to theoretically expected variables and not to the attitudinal dimension}\DIFdelend \DIFaddbegin \DIFadd{a prespecified set of expected relationships, including whether HPT dimensions relate differentially to external outcomes as theory predicts}\DIFaddend .
\end{enumerate}

The format itself may be a fairness mechanism. If scenario-based formats channel deliberative processing and override automatic attitude-retrieval \citep{Evans2013}, then the instrument format---not just the construct---buffers against motivated reasoning. However, this hypothesis has an important implication: if the format itself suppresses attitudinal responding, then the present fairness conclusion is conditional on the closed-ended Likert format and cannot be extended to open-ended HPT assessments where ideology could operate through different cognitive pathways.

Several design features may explain the instrument's resistance to
attitudinal contamination: the scenario provides contextual
information that channels deliberative processing, the Likert items
require evaluating propositions about a historical figure rather than
expressing personal attitudes, and the school assessment context may
activate an academic frame. These are hypotheses, not
conclusions---the present study was not designed to test them
experimentally, and doing so would require comparing structured and
open-ended formats using the same content.

An alternative explanation for the \DIFdelbegin \DIFdel{observed measurement equivalence
}\DIFdelend \DIFaddbegin \DIFadd{compatible invariance fit pattern
}\DIFaddend deserves consideration. In a school context where the scenario
involves a Nazi-adjacent historical figure, students may uniformly
avoid endorsing presentist responses---not because they are engaging
in genuine historical reasoning, but because \DIFdelbegin \DIFdel{contextualisation }\DIFdelend \DIFaddbegin \DIFadd{contextualization }\DIFaddend is the
socially safe response regardless of one's actual political
attitudes. If task-demand compliance produces equivalent response
patterns across ideology groups, the resulting \DIFdelbegin \DIFdel{measurement }\DIFdelend \DIFaddbegin \DIFadd{apparent }\DIFaddend invariance
would reflect shared social desirability pressure rather than genuine
absence of ideological influence on reasoning. The null SDR-5
correlations provide partial evidence against this interpretation, but
a brief self-report measure of social desirability may not capture the
specific task demands operating in a classroom assessment of
politically sensitive historical content.

\hypertarget{limitations}{%
\subsubsection{5.6 Limitations}\label{limitations}}

Several limitations qualify these findings.

\begin{enumerate}
\def\labelenumi{\arabic{enumi}.}

\item \emph{Single country and scenario.} \DIFdelbegin \DIFdel{The study uses one national
  context and one historical scenario ; whether the fairness findings
  generalise to other contexts remains to be established. }\DIFdelend \DIFaddbegin \DIFadd{Nine ratings of one Weimar
  scenario do not provide nine independent samples of the HPT content
  domain, so scenario-specific wording, knowledge, or moral salience
  cannot be separated from HPT. Future studies should use multiple
  scenarios spanning different periods, places, and political contexts.
}\DIFaddend

\item \emph{Restricted ideological range.} \DIFdelbegin \DIFdel{Only 17.7\% }\DIFdelend \DIFaddbegin \DIFadd{Depending on the measure,
  18.0\% (FR-LF), 37.5\% (KSA-3), or 25.1\% (their raw-score mean) }\DIFaddend of respondents
  scored above the scale midpoint\DIFdelbegin \DIFdel{on ideology measures}\DIFdelend , so the DIF and
  invariance analyses compare moderately elevated to low ideology, not
  genuine political extremism. The null ideology finding should be
  interpreted as ``no evidence of CIV in this range-restricted sample''
  rather than as a general demonstration that politically charged
  content is fair.

\item \emph{Limited DIF power.} With \emph{n} \(\approx\) 96 per
  group, \DIFdelbegin \DIFdel{the analysis was powered primarily to detect moderate-to-large
  effects (\(\Delta\beta\) \(\geq\) 0.5 logits); weak DIF effects
  cannot be confidently excluded}\DIFdelend \DIFaddbegin \DIFadd{DIF power was limited}\DIFaddend . A formal power simulation for the GRM
  likelihood-ratio test at Bonferroni-corrected \(\alpha\) was not
  conducted\DIFdelbegin \DIFdel{; future work should quantify }\DIFdelend \DIFaddbegin \DIFadd{, so }\DIFaddend the minimum detectable DIF magnitude \DIFdelbegin \DIFdel{for this design}\DIFdelend \DIFaddbegin \DIFadd{is unknown}\DIFaddend .

\item \emph{Gender imbalance.} The sample was 65.5\% male\DIFdelbegin \DIFdel{. Gender is
  a documented moderator of authoritarian attitude endorsement in
  adolescent samples}\DIFdelend , and no DIF
  or sensitivity analysis by gender was conducted. \DIFdelbegin \DIFdel{If male students drive both ideology variance and the null
  finding, the fairness
  conclusion may not hold for }\DIFdelend \DIFaddbegin \DIFadd{The fairness
  conclusion should therefore be replicated in a }\DIFaddend more gender-balanced
  sample\DIFdelbegin \DIFdel{s}\DIFdelend  .

\item \DIFaddbegin \emph{\DIFadd{Convenience sampling and unavailable socioeconomic data.}}
  \DIFadd{The sample over-represented upper-secondary students and was not
  probabilistic. Because no defensible socioeconomic indicator was
  collected, neither socioeconomic coverage nor selection along that
  dimension can be assessed.
}

\item \DIFaddend \emph{Modest reliability.} Subscale reliabilities (\(\alpha\) =
  .47--.64) constrain individual-level score interpretations; future
  adaptations should consider expanding subscale length. Omega
  hierarchical (\(\omega_h\)) from the bifactor model was not
  computed; this index would clarify what proportion of subscale
  variance is attributable to the general HPT factor versus specific
  subscale factors.

\item \emph{Method and format constraints.} All measures were
  self-report instruments administered in a single session. The
  closed-ended Likert format may not generalise to open-ended HPT
  assessments where ideology could operate through different cognitive
  mechanisms \citep{Smith2017Cognitive,Smith_et_al2018History}.

\item \emph{No response process evidence.} \DIFdelbegin \DIFdel{The critical unanswered
  question is whether Czech students actually invoke historical
  reasoning when responding---or whether the structured format produces
  contextualising responses through a different cognitive pathway.
  Cognitive interview studiescould resolve this
  .
}\DIFdelend \DIFaddbegin \DIFadd{Think-aloud or cognitive-
  interview studies, extending }\citet{Huijgen_et_al2017Toward}\DIFadd{, could
  compare the same content in rating and open-response formats. Evidence
  use and integration of conflicting historical considerations would
  support contextual reconstruction; explicit moral evaluation or
  affective engagement would identify processes outside this
  instrument's narrow cognitive scoring target; paraphrase and option-
  matching comments would indicate format-driven responding; and
  endorsement tied to one's present political position would indicate
  attitude congruence. These observable indicators could distinguish
  pathways that the present ratings cannot separate.
}\DIFaddend

\end{enumerate}

\hypertarget{recommendations}{%
\subsubsection{5.7 Recommendations}\label{recommendations}}

For researchers developing or adapting HPT instruments: (a) report \(\omega\) alongside \(\alpha\) for short subscales, with Spearman-Brown projections for context; (b) evaluate fairness through DIF and measurement invariance for theoretically relevant grouping variables; (c) consider revising weak-loading items (POP2, ROA2) while preserving the theoretical structure; (d) expand subscale length to improve reliability; and (e) replicate fairness evaluations across diverse samples, scenarios, and cultural contexts.

\hypertarget{conclusion}{%
\subsubsection{5.8 Conclusion}\label{conclusion}}

The Czech \DIFdelbegin \DIFdel{adaptation }\DIFdelend \DIFaddbegin \DIFadd{data were compatible with a correlated three-factor
representation }\DIFaddend of the Hartmann and Hasselhorn HPT \DIFdelbegin \DIFdel{instrument
preserves the three-factor structure of the original, produces }\DIFdelend \DIFaddbegin \DIFadd{item sets and produced }\DIFaddend scores
that relate to scenario-relevant historical knowledge in theoretically
expected ways\DIFdelbegin \DIFdel{, and functions with largely equivalent measurement
properties across students with different political attitudes}\DIFdelend . \DIFaddbegin \DIFadd{Fit-index patterns were compatible with measurement
invariance across ideology groups, but small extreme groups, unmodelled
clustering, and post-registration analysis prevent a firm equivalence
conclusion. }\DIFaddend At the
sensitivity levels afforded by the present sample and in the observed
ideological range, \DIFdelbegin \DIFdel{political ideology was unrelated to HPT scores
across all model specifications. }\DIFdelend \DIFaddbegin \DIFadd{right-authoritarian attitudes were unrelated to the
primary HPT composite. An unadjusted exploratory authoritarianism association
with one subscale was not subjected to the preregistered multiplicity
procedure.
}\DIFaddend

For researchers and assessment developers, the practical implication
is that structured scenario-based assessments can \DIFdelbegin \DIFdel{employ politically
sensitive historical content without detectable measurement bias---at least in mainstream classroom settings with restricted
ideological range}\DIFdelend \DIFaddbegin \DIFadd{be evaluated for
attitude-related measurement differences using complementary item-, model-,
and relationship-level checks. In this sample, those checks detected no bias
in the primary composite at the available sensitivity}\DIFaddend . Whether this \DIFdelbegin \DIFdel{conclusion }\DIFdelend \DIFaddbegin \DIFadd{pattern }\DIFaddend extends to more \DIFdelbegin \DIFdel{polarised
}\DIFdelend \DIFaddbegin \DIFadd{polarized
}\DIFaddend populations, open-ended assessment formats, or different historical
scenarios remains an open empirical question. The evaluation protocol
\DIFdelbegin \DIFdel{demonstrated }\DIFdelend \DIFaddbegin \DIFadd{illustrated }\DIFaddend here---CFA, DIF, MG-CFA, and \DIFdelbegin \DIFdel{convergent/discriminant
correlations---provides a replicable }\DIFdelend \DIFaddbegin \DIFadd{a prespecified network of
external relationships---provides a possible }\DIFaddend template for answering that
question wherever attitudinally charged content meets competence
assessment.

\hypertarget{data-availability-statement}{%
\subsection{Data Availability Statement}\label{data-availability-statement}}

De-identified data, annotated analysis scripts, the complete instrument battery (in Czech), supplementary materials, and all deviations from the original instrument are available on the Open Science Framework (https://doi.org/10.17605/OSF.IO/YNG37). Data are provided in machine-readable formats (RDS and xlsx) accompanied by a full variable codebook (PDF). The analysis pipeline consists of seven numbered R Markdown scripts that reproduce all reported results; a README file documents the directory structure, execution order, software requirements, and mapping of scripts to manuscript tables and figures. Supplementary materials (\DIFdelbegin \DIFdel{Table S1: GRM item parameters}\DIFdelend \DIFaddbegin \DIFadd{Tables S1--S5}\DIFaddend ) are included in the OSF repository. The study design, hypotheses, and analysis plan were preregistered on OSF prior to data collection\DIFaddbegin \DIFadd{; deviations and omissions are disclosed in Table S5}\DIFaddend .

\hypertarget{conflict-of-interest-disclosure}{%
\subsection{Conflict of Interest Disclosure}\label{conflict-of-interest-disclosure}}

The author declares that he complies with the PCI rule of having no financial conflicts of interest in relation to the content of the article, and has no non-financial conflicts of interest.

\hypertarget{ethics-statement}{%
\subsection{Ethics Statement}\label{ethics-statement}}

This study involved secondary school students completing questionnaires during regular class periods. Under Czech law (Act No.~110/2019 Sb., on the processing of personal data), ethics committee approval is not required for non-interventional educational research involving anonymous questionnaires where no personal data are collected. Participation was voluntary; students were informed about the study's purpose and their right to decline participation without consequences. All responses were de-identified at the point of data entry. No personally identifiable information was collected or stored. The study followed the principles of the Declaration of Helsinki for research involving human participants.

\hypertarget{funding}{%
\subsection{Funding}\label{funding}}

This research received no specific grant from any funding agency in the public, commercial, or not-for-profit sectors.

\hypertarget{author-contributions-credit}{%
\subsection{Author Contributions (CRediT)}\label{author-contributions-credit}}

Juda Kaleta: Conceptualization, Methodology, Software, Formal analysis, Investigation, Data curation, Writing --- original draft, Writing --- review \& editing, Visualization, Project administration.

\hypertarget{acknowledgements}{%
\subsection{Acknowledgements}\label{acknowledgements}}

The author thanks the participating schools, teachers, and students for their cooperation. The translation panel members and expert reviewers who contributed to the adaptation process are acknowledged with gratitude.

\hypertarget{ai-use-disclosure}{%
\subsection{AI Use Disclosure}\label{ai-use-disclosure}}

The following AI tools were used during the preparation of this manuscript:

\begin{itemize}
\tightlist
\item
  \emph{Brainstorming and scripting assistance:} Claude (Anthropic) ---
  used for exploring conceptual framing, structural alternatives, and
  scripting support for data processing and analysis workflows.
\item
  \emph{Language editing:} Grammarly --- used for grammar and style
  checking.
\end{itemize}

All research design decisions, analytical choices, theoretical interpretations, and the final text are solely the author's responsibility. The author reviewed, revised, and approved all content.

\bibliographystyle{apacite}
\bibliography{bibliography,extras}

\newpage

\hypertarget{top-checklist}{%
\section*{TOP Checklist and Disclosures Form}\label{top-checklist}}

This checklist must be completed and included as an appendix at the end of your preprint prior to submission to PCI Psychology. Manuscripts without this document included as an appendix will be returned to authors without review.

The policy of PCI Psychology is to recommend papers only if the data, methods used in the analysis, and any digital materials used to conduct the research are clearly and precisely documented and are maximally available to any researcher for purposes of reproducing the results or replicating the procedure. PCI Psych follows the principle of ``as open as possible, as closed as necessary.''

\bigskip

\textbf{First author name (last/family, first/given):} Kaleta, Juda

\textbf{Preprint DOI or URL:} \url{https://doi.org/10.17605/OSF.IO/YNG37}

\bigskip

\hypertarget{section-1-data}{%
\subsection*{Section 1: Data}\label{section-1-data}}

\textbf{Does your manuscript contain reports of any data?}

\textbf{[X] Yes} (continue with next question)

\quad $\square$ No (skip to Section 2)

\bigskip

\textbf{Are appropriately \DIFdelbegin \DIFdel{anonymised }\DIFdelend \DIFaddbegin \DIFadd{anonymized }\DIFaddend raw data available within a trusted digital repository?}

\textbf{[X] Yes, available at this link:} \url{https://doi.org/10.17605/OSF.IO/YNG37}

\quad $\square$ No, justification:

\bigskip

\textbf{Are third-party data cited in the manuscript, with a DOI?} (e.g., for preexisting data, data deposited in a repository)

\quad $\square$ Yes, the DOI is as follows:

\textbf{[X] No, justification:} The data were collected specifically for this study; no third-party datasets were used.

\bigskip

\textbf{Is there a data dictionary and/or readme file included with the data to make it interpretable?}

\textbf{[X] Yes, available at this link:} \url{https://doi.org/10.17605/OSF.IO/YNG37} (full variable codebook in PDF and README file)

\quad $\square$ No, justification:

\bigskip

\textbf{Do you indicate in the manuscript how the sample size was determined?}

\textbf{[X] Yes.} (Section 3.4: a priori power analysis targeting 80\% power to detect $\beta \geq .15$)

\quad $\square$ No, justification:

\bigskip

\textbf{Do you report all data exclusions (e.g., outliers, careless responders)?}

\textbf{[X] Yes.} (Section 3.4 and \DIFdelbegin \DIFdel{4.6: }\DIFdelend \DIFaddbegin \DIFadd{Table S5: the }\DIFaddend middle ideology tertile \DIFdelbegin \DIFdel{excluded for DIF extreme-group design; all exclusions reported}\DIFdelend \DIFaddbegin \DIFadd{was excluded for post-registration DIF and MG-CFA analyses; this analysis-specific exclusion is reported.}\DIFaddend )

\quad $\square$ No, justification:

\bigskip

\textbf{Do you report all inclusion/exclusion criteria and when they were established?}

\textbf{[X] Yes.} (\DIFdelbegin \DIFdel{Inclusion/exclusion criteria }\DIFdelend \DIFaddbegin \DIFadd{Participant eligibility is described in Section 3.2 and was preregistered; post-registration analysis-specific exclusions }\DIFaddend are described in Section\DIFdelbegin \DIFdel{s}\DIFdelend  3.4 \DIFdelbegin \DIFdel{--3.5 and were preregistered prior to data collection}\DIFdelend \DIFaddbegin \DIFadd{and Table S5.}\DIFaddend )

\quad $\square$ No, justification:

\bigskip

\textbf{Are all measures, questions, and/or conditions used in the study described in the manuscript or available in the supplemental material?}

\textbf{[X] Yes.} (All instruments described in Sections 3.1--3.3; full Czech-language instrument battery available on OSF: \url{https://doi.org/10.17605/OSF.IO/YNG37})

\quad $\square$ No, justification:

\bigskip

\hypertarget{section-2-analysis}{%
\subsection*{Section 2: Analysis Scripts/Code/Codebooks}\label{section-2-analysis}}

\textbf{Does your manuscript contain any analysis of quantitative or qualitative data?}

\textbf{[X] Yes} (continue with next question)

\quad $\square$ No (skip to Section 3)

\bigskip

\textbf{Are third-party analysis scripts/code (e.g., R, Stata), codebooks, or other relevant documentation available within a trusted digital repository?}

\textbf{[X] Yes, available at this link:} \url{https://doi.org/10.17605/OSF.IO/YNG37} (7 R Markdown pipeline scripts, 4 figure scripts, 2 supplementary analysis scripts, and a Makefile)

\quad $\square$ No, justification:

\bigskip

\textbf{Are the analysis scripts/code (e.g., R, Stata), codebooks, or other relevant documentation cited in the manuscript, with a DOI?}

\textbf{[X] Yes, the DOI is as follows:} \url{https://doi.org/10.17605/OSF.IO/YNG37}

\quad $\square$ No, justification:

\bigskip

\hypertarget{section-3-materials}{%
\subsection*{Section 3: Study Materials}\label{section-3-materials}}

\textbf{Does your manuscript contain any research materials (e.g., stimuli, programming code, questionnaires, interview protocols)?}

\textbf{[X] Yes} (continue with next question)

\quad $\square$ No (skip to Section 4)

\bigskip

\textbf{Are all study materials and descriptions of study procedures available within a trusted digital repository?}

\textbf{[X] Yes, available at this link:} \url{https://doi.org/10.17605/OSF.IO/YNG37} (Czech-language versions of all questionnaires: HPT, FR-LF, KSA-3, SDR-5, and historical knowledge test)

\quad $\square$ No, justification:

\bigskip

\textbf{Are all third-party study materials, descriptions of study procedures, or other relevant documents cited in the manuscript, with a DOI?}

\quad $\square$ Yes, the DOI is as follows:

\textbf{[X] No, justification:} Third-party instruments \DIFdelbegin \DIFdel{(HPT, FR-LF, KSA-3, SDR-5) }\DIFdelend are cited via standard bibliography references (Hartmann \& Hasselhorn, 2008; \DIFdelbegin \DIFdel{Becker }\DIFdelend \DIFaddbegin \DIFadd{Decker }\DIFaddend et al., \DIFdelbegin \DIFdel{2020; Lederer, 2002; Stöber, 2001}\DIFdelend \DIFaddbegin \DIFadd{2013; Beierlein et al., 2014; Hays et al., 1989}\DIFaddend ); DOIs are included where available. The adapted Czech versions are deposited on OSF (\url{https://doi.org/10.17605/OSF.IO/YNG37}).

\bigskip

\hypertarget{section-4-preregistration}{%
\subsection*{Section 4: Preregistration}\label{section-4-preregistration}}

\textbf{Were any aspects of your manuscript preregistered?}

\textbf{[X] Yes} (continue with next question)

\quad $\square$ No (do not complete the rest of the form)

\bigskip

\textbf{Does the manuscript contain an accessible link to the preregistration?}

\textbf{[X] Yes, available at this link:} \url{https://osf.io/zsngy}

\quad $\square$ No, justification:

\bigskip

\textbf{Do you clearly indicate in the manuscript which parts were preregistered and which parts were not?}

\textbf{[X] Yes.} (\DIFdelbegin \DIFdel{Preregistered components include hypotheses H1--H2, primary multilevel models, DIF/MG-CFA analyses, and sample size determination; the TOST equivalence analysis is explicitly labelled ``exploratory'' in Section 4.10 as it used a post-hoc SESOI wider than the preregistered power target}\DIFdelend \DIFaddbegin \DIFadd{Section 2.5 gives a concise disclosure of post-registration changes and omitted registered analyses. Table S5 distinguishes registered tests, omissions, implemented modifications, and post-registration analyses.}\DIFaddend )

\quad $\square$ No, justification:

\bigskip

\textbf{Are all preregistered analyses reported in the text or linked in the supplemental material?}

\DIFaddbegin \DIFadd{\quad $\square$ Yes.
}

\DIFaddend \textbf{[X] \DIFdelbegin \DIFdel{Yes.}\DIFdelend \DIFaddbegin \DIFadd{No, justification:}\DIFaddend } \DIFdelbegin \DIFdel{(All preregistered hypotheses and analyses are reported in Sections 4.1--4.10 and the supplementary materials)
}\DIFdelend \DIFaddbegin \DIFadd{The registered continuous-ideology MIMIC model, complete moderation and simple-slope analyses, clustered correlation comparison, complete covariate set, one-tailed tests, Benjamini--Hochberg correction, and registered missing-data rule were not all implemented. Table S5 reports these omissions and the analyses used instead.
}\DIFaddend

\DIFdelbegin \DIFdel{\quad $\square$ No, justification:
}%DIFDELCMD <

%DIFDELCMD < %%%
\DIFdelend \bigskip

\textbf{Are all deviations from the preregistration plan clearly disclosed in the manuscript (either in text or in a table)?}

\textbf{[X] Yes.} (\DIFdelbegin \DIFdel{One deviation: TOST SESOI set post hoc at }\DIFdelend \DIFaddbegin \DIFadd{Section 2.5 and Table S5 disclose the changed inferential hierarchy, omitted registered procedures, instrument discrepancies, and post-registration additions. TOST and its }\DIFaddend $\pm$0.20 \DIFdelbegin \DIFdel{instead of the preregistered power target of $\beta \geq .15$; disclosed in Section 4.10 and the analysis is labelled exploratory}\DIFdelend \DIFaddbegin \DIFadd{SESOI are explicitly labelled exploratory.}\DIFaddend )

\end{document}
